\documentclass[11pt,a4paper,oneside]{report}
\usepackage[utf8]{inputenc}
\usepackage[vietnamese]{babel}
\usepackage[T5]{fontenc}
\usepackage{geometry}
\geometry{a4paper, top=2.5cm, bottom=2.5cm, left=2.5cm, right=2.5cm}
\usepackage{graphicx}
\usepackage{booktabs}
\usepackage{tabularx}
\usepackage{array}
\usepackage{xcolor}
\usepackage{tcolorbox}
\tcbuselibrary{skins,breakable}
\usepackage{enumitem}
\usepackage{amsmath,amssymb}
\usepackage{fancyhdr}
\usepackage{titlesec}
\usepackage{hyperref}

\hypersetup{
    colorlinks=true,
    linkcolor=blue!75!black,
    filecolor=magenta,
    urlcolor=teal,
    citecolor=green!60!black
}

\newtcolorbox{myquote}[1][]{
  colback=blue!3!white,
  colframe=blue!60!black,
  boxrule=0.8pt,
  arc=2mm,
  left=10pt,
  right=10pt,
  top=8pt,
  bottom=8pt,
  breakable,
  #1
}

\pagestyle{fancy}
\fancyhf{}
\fancyhead[L]{\small\textsf{Thiết Kế Giao Diện Người Dùng Đồ Họa (UI/UX - GUID)}}
\fancyhead[R]{\small\thepage}
\fancyfoot[C]{\small\textsf{Tài liệu học tập và ôn thi HK9}}
\renewcommand{\headrulewidth}{0.4pt}
\renewcommand{\footrulewidth}{0.4pt}

\titleformat{\chapter}[display]
  {\normalfont\bfseries\color{blue!80!black}}
  {\Large\MakeUppercase{\chaptertitlename} \thechapter}{1ex}
  {\titlerule\vspace{1ex}\huge}
  [\vspace{1ex}\titlerule]

\titlespacing*{\chapter}{0pt}{-20pt}{20pt}

\begin{document}

\begin{titlepage}
    \centering
    \vspace*{1.5cm}
    {\scshape\LARGE TRƯỜNG ĐẠI HỌC KỸ THUẬT - CÔNG NGHỆ CẦN THƠ \par}
    {\scshape\large KHOA CÔNG NGHỆ THÔNG TIN \par}
    \vspace{1.5cm}
    \hrule height 1.5pt
    \vspace{0.6cm}
    {\huge\bfseries GIÁO TRÌNH HỌC TẬP VÀ ÔN THI TOÀN DIỆN\par}
    \vspace{0.4cm}
    {\LARGE\bfseries THIẾT KẾ GIAO DIỆN NGƯỜI DÙNG ĐỒ HỌA\par}
    {\Large\textit{(Graphical User Interface Design - GUID / UI-UX Design)}\par}
    \vspace{0.6cm}
    \hrule height 1.5pt
    \vspace{2cm}
    
    \begin{minipage}{0.85\textwidth}
        \large
        \textbf{Môn học:} Thiết kế giao diện người dùng đồ họa (GUID)\\
        \textbf{Khóa học:} HK9 - Công nghệ thông tin\\
        \textbf{Nội dung:} 4 Chương lý thuyết, Hệ thống 10 nguyên tắc Heuristic, Case Studies thực tế, Câu hỏi tình huống có đáp án và Bài tập lớn.\\
        \textbf{Tài liệu tham khảo:} Bài giảng chính thức của giảng viên bộ môn, Steve Krug, Jakob Nielsen, Ben Shneiderman.
    \end{minipage}
    
    \vfill
    {\large Cần Thơ, Năm 2026 \par}
\end{titlepage}

\tableofcontents
\newpage

\chapter{CHƯƠNG 1: GIỚI THIỆU VỀ UI/UX (INTRODUCTION TO UI/UX)}

\section{TỔNG QUAN MÔN HỌC}
\begin{itemize}[leftmargin=1.5em]
\item \textbf{Môn học:} Thiết kế giao diện người dùng đồ họa (\textit{Graphical User Interface Design - GUID})  
\item \textbf{Khái niệm cốt lõi:} GUI - Graphical User Interface  
\end{itemize}

\begin{myquote}
\textbf{\large Bốn mục tiêu của giáo dục theo UNESCO:}\\
1. \textbf{Học để biết (Learning to know):} Nắm vững các kiến thức, nguyên lý và quy chuẩn thiết kế.
2. \textbf{Học để làm (Learning to do):} Ứng dụng thực hành thành thạo các công cụ và quy trình thiết kế giao diện thực tế.
3. \textbf{Học để chung sống (Learning to live together):} Làm việc nhóm hiệu quả, thấu hiểu và đồng cảm sâu sắc với người dùng.
4. \textbf{Học để tự khẳng định mình (Learning to be):} Phát triển tư duy thẩm mỹ, sáng tạo và bản sắc thiết kế riêng biệt.
\end{myquote}

\bigskip\hrule\bigskip

\section{1.1 Khái niệm về sản phẩm và 3 mục tiêu cốt lõi}

\subsection{1.1.1 Onboarding \& Case Studies thực tế}

\subsubsection{Case Study 1: Câu chuyện gia nhập thị trường của BAEMIN Việt Nam}
\begin{itemize}[leftmargin=1.5em]
\item \textbf{Bối cảnh:} BAEMIN bước chân vào thị trường giao nhận đồ ăn (food delivery) tại Việt Nam khi nơi đây đã có sự thống trị của những "gã khổng lồ" như GrabFood, Now Delivery (sau này là ShopeeFood).
\item \textbf{Chiến thuật tiếp cận độc đáo:}
\item Thay vì dàn trải nguồn lực, các nhà lãnh đạo BAEMIN khi vào thị trường đầu tiên tại Việt Nam (TP. Hồ Chí Minh) đã chọn chiến thuật \textbf{"tấn công riêng lẻ"}, tập trung nguồn lực đánh chiếm từng quận cụ thể (bắt đầu từ Quận 1, Quận 3).
\item Triết lý cạnh tranh: \textbf{Chiến lược "Đứng nhất từng quận"}, thay vì đứng thứ hai hay thứ ba ở tất cả các quận.
\item Sau khi đạt được các tín hiệu khả quan, có tập dữ liệu vững chắc về hành vi khách hàng và thị trường, BAEMIN mới mở rộng ra các quận lân cận và tiến tới các thành phố tiếp theo như Thủ đô Hà Nội, TP. Đà Nẵng.
\item \textbf{Các nền tảng dịch vụ giao đồ ăn tiêu biểu:} GrabFood, Now, BAEMIN, GO-FOOD, Lozi.
\end{itemize}

\begin{figure}[htbp]
\centering
\includegraphics[width=0.75\textwidth]{images/image1.png}
\caption{Ứng dụng BAEMIN và chiến lược chinh phục người dùng}
\end{figure}

\subsubsection{Case Study 2: Bài học từ mạng xã hội LOTUS}
\begin{itemize}[leftmargin=1.5em]
\item \textbf{Bối cảnh:} Lotus ra mắt với sứ mệnh xây dựng mạng xã hội "Made in Vietnam" đầy tham vọng, được hỗ trợ bởi nền tảng điện toán đám mây BizFly Cloud mạnh mẽ.
\item \textbf{Quy mô đầu tư:} Được nghiên cứu và phát triển suốt hơn một năm bởi đội ngũ hơn 200 kỹ sư công nghệ hàng đầu trong các lĩnh vực Di động (Mobile), Trí tuệ nhân tạo (AI), Dữ liệu lớn (Big Data), Điện toán đám mây (Cloud Computing).
\item \textbf{Chiến dịch ra mắt:} Ra mắt vô cùng rầm rộ với phương châm cốt lõi: \textit{"NỘI DUNG LÀ VUA"}.
\item \textbf{Vì sao Lotus thất bại tại Việt Nam?}
\item Mặc dù có tiềm lực công nghệ và tài chính lớn, sản phẩm lại thiếu sự thấu hiểu sâu sắc về trải nghiệm người dùng (UX) thực tế.
\item Giao diện phức tạp, luồng người dùng rườm rà, thiếu tính năng giữ chân khách hàng tự nhiên.
\item Thất bại của Lotus minh chứng cho quy luật sống còn: Công nghệ hiện đại không thể cứu vãn một trải nghiệm người dùng kém và thiếu giá trị thực tiễn.
\end{itemize}

\bigskip\hrule\bigskip

\subsection{1.1.2 Sản phẩm là gì? Ba mục tiêu cốt lõi}
\begin{itemize}[leftmargin=1.5em]
\item \textbf{Định nghĩa sản phẩm (Product):} Sản phẩm sinh ra là để đáp ứng, giải quyết những nhu cầu, vấn đề cụ thể nào đó của con người.
\item \textbf{Ba mục tiêu cốt lõi của một sản phẩm thành công:}
\end{itemize}
\begin{enumerate}[leftmargin=1.5em]
\item \textbf{Mục tiêu Trải nghiệm người dùng (User Experience - UX):} Sản phẩm có giải quyết đúng vấn đề của người dùng một cách dễ dàng, thuận tiện và mang lại cảm xúc tích cực hay không?
\item \textbf{Mục tiêu Kinh doanh (Business Goals):} Sản phẩm có mang lại lợi nhuận, doanh thu hoặc giá trị chiến lược cho doanh nghiệp hay không?
\item \textbf{Mục tiêu Khả thi công nghệ (Technology / Feasibility):} Năng lực kỹ thuật, hạ tầng phần cứng và phần mềm có đủ khả năng xây dựng, vận hành và mở rộng sản phẩm ổn định hay không?
\end{enumerate}

\begin{figure}[htbp]
\centering
\includegraphics[width=0.75\textwidth]{images/image2.png}
\caption{Ba mục tiêu cốt lõi của sản phẩm}
\end{figure}

\subsubsection{Phân tích hệ quả khi thiếu một trong ba yếu tố:}
\begin{itemize}[leftmargin=1.5em]
\item \textbf{Thiếu "User" (Người dùng):} Sản phẩm áp dụng công nghệ rất cao siêu, mô hình kinh doanh trên giấy rất hoàn hảo nhưng người dùng không có nhu cầu hoặc thấy quá khó dùng $
\end{itemize}
ightarrow$ Sản phẩm chết yểu (ví dụ điển hình như Lotus).

\begin{itemize}[leftmargin=1.5em]
\item \textbf{Thiếu "Business" (Kinh doanh):} Người dùng rất thích, công nghệ rất mượt mà nhưng không có dòng tiền, chi phí vận hành quá lớn không thể bù đắp $
\end{itemize}
ightarrow$ Doanh nghiệp cạn vốn và phá sản.

\begin{itemize}[leftmargin=1.5em]
\item \textbf{Thiếu "Technology" (Công nghệ):} Nhu cầu người dùng rất lớn, cơ hội kinh doanh khổng lồ nhưng đội ngũ không thể hiện thực hóa về mặt kỹ thuật, hệ thống giật lag, sập server liên tục $
\end{itemize}
ightarrow$ Đánh mất niềm tin và khách hàng rời bỏ.


\bigskip\hrule\bigskip

\section{1.2 Sản phẩm phần mềm và các yếu tố phát triển bền vững}

\subsection{1.2.1 Sản phẩm phần mềm là gì?}
\textbf{Phần mềm (Software)} là tập hợp các dữ liệu, chương trình hoặc các câu lệnh hướng dẫn máy tính, giúp cho các thiết bị phần cứng biết cách vận hành và xử lý công việc.


\subsection{1.2.2 Mười yếu tố quan trọng giúp công ty công nghệ duy trì sự phát triển bền vững:}
\begin{enumerate}[leftmargin=1.5em]
\item \textbf{Nghiên cứu và Phát triển (R\&D):}
\end{enumerate}
\begin{itemize}[leftmargin=1.5em]
\item Không ngừng đổi mới sáng tạo sản phẩm.
\item Chủ động nghiên cứu đón đầu các công nghệ mới nổi.
\end{itemize}
\begin{enumerate}[leftmargin=1.5em]
\item \textbf{Tập trung vào Khách hàng (Customer-Centric):}
\end{enumerate}
\begin{itemize}[leftmargin=1.5em]
\item Thấu hiểu sâu sắc nhu cầu thực và "nỗi đau" (pain points) của khách hàng.
\item Cung cấp dịch vụ chăm sóc và hỗ trợ khách hàng xuất sắc.
\end{itemize}
\begin{enumerate}[leftmargin=1.5em]
\item \textbf{Phát triển Nguồn nhân lực:}
\end{enumerate}
\begin{itemize}[leftmargin=1.5em]
\item Chú trọng đào tạo chuyên môn và nâng cao kỹ năng liên tục.
\item Xây dựng văn hóa và môi trường làm việc tích cực, sáng tạo.
\end{itemize}
\begin{enumerate}[leftmargin=1.5em]
\item \textbf{Quản lý Tài chính hiệu quả:}
\end{enumerate}
\begin{itemize}[leftmargin=1.5em]
\item Kiểm soát chặt chẽ dòng tiền và chi phí vận hành.
\item Đầu tư chiến lược có trọng tâm, tối ưu hóa tỷ suất sinh lời (ROI).
\end{itemize}
\begin{enumerate}[leftmargin=1.5em]
\item \textbf{Chiến lược Tiếp thị và Thương hiệu:}
\end{enumerate}
\begin{itemize}[leftmargin=1.5em]
\item Định vị và xây dựng thương hiệu mạnh mẽ, uy tín.
\item Triển khai chiến lược tiếp thị số (Digital Marketing) thông minh và linh hoạt.
\end{itemize}
\begin{enumerate}[leftmargin=1.5em]
\item \textbf{Đa dạng hóa Sản phẩm và Dịch vụ:}
\end{enumerate}
\begin{itemize}[leftmargin=1.5em]
\item Mở rộng danh mục sản phẩm bổ trợ lẫn nhau trong hệ sinh thái.
\item Khai phá và thâm nhập các phân khúc thị trường mới.
\end{itemize}
\begin{enumerate}[leftmargin=1.5em]
\item \textbf{Hợp tác và Liên minh chiến lược:}
\end{enumerate}
\begin{itemize}[leftmargin=1.5em]
\item Xây dựng liên minh vững chắc với các đối tác chiến lược.
\item Tận dụng các cơ hội M\&A (Mua bán \& Sáp nhập) để gia tăng sức mạnh công nghệ và thị phần.
\end{itemize}
\begin{enumerate}[leftmargin=1.5em]
\item \textbf{Thích ứng với Thay đổi:}
\end{enumerate}
\begin{itemize}[leftmargin=1.5em]
\item Luôn duy trì sự linh hoạt và khả năng thích nghi cao trước biến động thị trường.
\item Đẩy mạnh ứng dụng công nghệ số và chuyển đổi mô hình số hóa.
\end{itemize}
\begin{enumerate}[leftmargin=1.5em]
\item \textbf{Đạo đức và Trách nhiệm xã hội (CSR):}
\end{enumerate}
\begin{itemize}[leftmargin=1.5em]
\item Nghiêm túc tuân thủ pháp luật, bảo vệ dữ liệu và quyền riêng tư của người dùng.
\item Đóng góp tích cực cho cộng đồng, thúc đẩy sự phát triển bền vững.
\end{itemize}
\begin{enumerate}[leftmargin=1.5em]
\item \textbf{Chiến lược Dài hạn:}
\end{enumerate}
\begin{itemize}[leftmargin=1.5em]
\item Xác lập tầm nhìn (Vision) và sứ mệnh (Mission) rõ ràng.
\item Lập kế hoạch hành động chiến lược theo từng giai đoạn cụ thể.
\end{itemize}

\bigskip\hrule\bigskip

\section{1.3 Phân loại các danh mục sản phẩm và dịch vụ công nghệ}

Các công ty công nghệ hiện đại kiến tạo nên một hệ sinh thái sản phẩm và dịch vụ đa dạng bao gồm:

\begin{enumerate}[leftmargin=1.5em]
\item \textbf{Phần mềm (Software):}
\end{enumerate}
\begin{itemize}[leftmargin=1.5em]
\item \textit{Ứng dụng di động (Mobile Apps):} Chạy trên smartphone và máy tính bảng (Messenger, WhatsApp, TikTok, Grab...).
\item \textit{Phần mềm máy tính (Desktop Software):} Cài đặt trên máy tính như Hệ điều hành (Windows, macOS, Linux), phần mềm văn phòng, đồ họa (Photoshop, AutoCAD).
\item \textit{Ứng dụng trên nền tảng đám mây (Cloud Apps / SaaS):} Dịch vụ trực tuyến lưu trữ và xử lý trên mây (Google Drive, Dropbox, Notion).
\end{itemize}
\begin{enumerate}[leftmargin=1.5em]
\item \textbf{Thiết bị phần cứng (Hardware Devices):}
\end{enumerate}
\begin{itemize}[leftmargin=1.5em]
\item \textit{Máy tính \& Máy tính bảng:} MacBook, Surface, ThinkPad, iPad...
\item \textit{Điện thoại thông minh (Smartphones):} iPhone, Samsung Galaxy, Google Pixel...
\item \textit{Thiết bị đeo thông minh (Wearables):} Apple Watch, Galaxy Watch, thiết bị theo dõi sức khỏe Garmin.
\end{itemize}
\begin{enumerate}[leftmargin=1.5em]
\item \textbf{Dịch vụ điện toán đám mây (Cloud Services):}
\end{enumerate}
\begin{itemize}[leftmargin=1.5em]
\item IaaS, PaaS, SaaS do các tập đoàn lớn cung cấp: Amazon Web Services (AWS), Microsoft Azure, Google Cloud Platform (GCP).
\end{itemize}
\begin{enumerate}[leftmargin=1.5em]
\item \textbf{Trí tuệ nhân tạo (AI) và Học máy (Machine Learning):}
\end{enumerate}
\begin{itemize}[leftmargin=1.5em]
\item Trợ lý ảo thông minh: Siri, Google Assistant, Amazon Alexa.
\item Các mô hình tạo sinh (Generative AI), dịch thuật tự động, phân tích dữ liệu lớn dự đoán xu hướng.
\end{itemize}
\begin{enumerate}[leftmargin=1.5em]
\item \textbf{Nền tảng Thương mại điện tử (E-commerce Platforms):}
\end{enumerate}
\begin{itemize}[leftmargin=1.5em]
\item Các sàn thương mại điện tử: Amazon, Taobao, Shopee, Lazada, Tiki.
\item Hệ thống thanh toán trực tuyến: PayPal, Apple Pay, MoMo, VNPay.
\end{itemize}
\begin{enumerate}[leftmargin=1.5em]
\item \textbf{Truyền thông xã hội và Giải trí số:}
\end{enumerate}
\begin{itemize}[leftmargin=1.5em]
\item Mạng xã hội: Facebook, Instagram, LinkedIn, X (Twitter).
\item Dịch vụ truyền phát trực tuyến (Streaming): Netflix, Spotify, YouTube.
\end{itemize}
\begin{enumerate}[leftmargin=1.5em]
\item \textbf{Công nghệ Di chuyển \& Phương tiện tự hành:}
\end{enumerate}
\begin{itemize}[leftmargin=1.5em]
\item Xe điện thông minh (Tesla, Rivian, VinFast).
\item Công nghệ xe tự hành ứng dụng cảm biến LiDAR, thị giác máy tính và AI (Waymo, Cruise).
\end{itemize}
\begin{enumerate}[leftmargin=1.5em]
\item \textbf{Robotics và Tự động hóa:}
\end{enumerate}
\begin{itemize}[leftmargin=1.5em]
\item Robot công nghiệp: Tự động hóa dây chuyền lắp ráp nhà máy.
\item Robot dịch vụ: Robot phẫu thuật y tế, robot phục vụ nhà hàng, robot hút bụi thông minh.
\end{itemize}
\begin{enumerate}[leftmargin=1.5em]
\item \textbf{An ninh mạng (Cybersecurity):}
\end{enumerate}
\begin{itemize}[leftmargin=1.5em]
\item Phần mềm bảo mật, diệt virus và tường lửa: Symantec, Kaspersky, McAfee.
\item Dịch vụ bảo mật dữ liệu, giám sát an toàn thông tin doanh nghiệp (SOC).
\end{itemize}
\begin{enumerate}[leftmargin=1.5em]
\item \textbf{Nền tảng Giao tiếp và Làm việc từ xa:}
\end{enumerate}
\begin{itemize}[leftmargin=1.5em]
\item Ứng dụng họp trực tuyến: Zoom, Microsoft Teams, Google Meet.
\item Công cụ quản lý dự án và luồng công việc: Jira, Asana, Trello, Slack.
\end{itemize}

\bigskip\hrule\bigskip

\section{1.4 Vai trò và nhiệm vụ của UI/UX Designer trong công ty công nghệ}

Trong một doanh nghiệp công nghệ, chuyên viên thiết kế UI/UX đảm nhận 7 nhóm trách nhiệm cốt lõi:

\begin{enumerate}[leftmargin=1.5em]
\item \textbf{Nghiên cứu người dùng (User Research):}
\end{enumerate}
\begin{itemize}[leftmargin=1.5em]
\item Phân tích nhu cầu, hành vi, bối cảnh sử dụng của nhóm người dùng mục tiêu.
\item Xây dựng các bản mô tả chân dung khách hàng điển hình (\textbf{User Persona}) để định hướng thiết kế.
\end{itemize}
\begin{enumerate}[leftmargin=1.5em]
\item \textbf{Thiết kế trải nghiệm người dùng (UX Design):}
\end{enumerate}
\begin{itemize}[leftmargin=1.5em]
\item Thiết lập luồng thao tác của người dùng (\textbf{User Flow / User Journey}).
\item Xây dựng khung xương định vị bố cục (\textbf{Wireframe}) từ cấp độ thấp đến cao.
\item Thiết kế nguyên mẫu thử nghiệm tương tác (\textbf{Interactive Prototyping}).
\end{itemize}
\begin{enumerate}[leftmargin=1.5em]
\item \textbf{Thiết kế giao diện người dùng (UI Design):}
\end{enumerate}
\begin{itemize}[leftmargin=1.5em]
\item Thiết kế trực quan: Lựa chọn phong cách thẩm mỹ, hệ màu (Color Palette), kiểu chữ (Typography), hệ thống Icon và hình ảnh.
\item Xây dựng hệ thống thiết kế (\textbf{Design System}) chuẩn mực nhằm đảm bảo sự nhất quán trên toàn bộ sản phẩm.
\item Tinh chỉnh và tối ưu hóa giao diện hiển thị sắc nét trên đa dạng độ phân giải.
\end{itemize}
\begin{enumerate}[leftmargin=1.5em]
\item \textbf{Kiểm thử và đánh giá (Testing \& Evaluation):}
\end{enumerate}
\begin{itemize}[leftmargin=1.5em]
\item Thực hiện kiểm thử tính khả dụng (\textbf{Usability Testing}) với người dùng thực tế.
\item Thu thập, phân tích dữ liệu định tính và định lượng từ phản hồi của người dùng để cải tiến sản phẩm.
\end{itemize}
\begin{enumerate}[leftmargin=1.5em]
\item \textbf{Cộng tác liên chức năng (Cross-functional Collaboration):}
\end{enumerate}
\begin{itemize}[leftmargin=1.5em]
\item Làm việc chặt chẽ với đội ngũ Lập trình viên (Developers) để hiện thực hóa bản vẽ đúng chuẩn kỹ thuật.
\item Trình bày, thuyết phục và bảo vệ ý tưởng thiết kế trước các bên liên quan (Product Manager, Stakeholders).
\end{itemize}
\begin{enumerate}[leftmargin=1.5em]
\item \textbf{Cập nhật xu hướng và công nghệ:}
\end{enumerate}
\begin{itemize}[leftmargin=1.5em]
\item Liên tục theo dõi các xu hướng thiết kế giao diện và tương tác hiện đại trên thế giới.
\item Nâng cao trình độ làm chủ các bộ công cụ thiết kế chuyên nghiệp (Figma, Sketch, Adobe XD...).
\end{itemize}
\begin{enumerate}[leftmargin=1.5em]
\item \textbf{Quản lý dự án và tiến độ:}
\end{enumerate}
\begin{itemize}[leftmargin=1.5em]
\item Lập kế hoạch phân bổ thời gian, đảm bảo hoàn thành nhiệm vụ đúng tiến độ cam kết.
\item Linh hoạt xử lý các yêu cầu thay đổi nghiệp vụ phát sinh trong quá trình triển khai dự án.
\end{itemize}

\bigskip\hrule\bigskip

\section{1.5 Các công cụ hỗ trợ thiết kế giao diện đồ họa (GUI Tools)}

Có nhiều phần mềm chuyên dụng hỗ trợ thiết kế UI/UX từ giai đoạn sơ thảo đến hoàn thiện:

\begin{itemize}[leftmargin=1.5em]
\item \textbf{Figma:} Công cụ thiết kế giao diện dựa trên nền tảng đám mây phổ biến nhất hiện nay, hỗ trợ làm việc cộng tác thời gian thực vượt trội.
\item \textbf{Sketch:} Ứng dụng thiết kế giao diện chuẩn mực chuyên biệt trên hệ điều hành macOS.
\item \textbf{Adobe XD:} Bộ công cụ thiết kế vector và tạo prototype tương tác của Adobe.
\item \textbf{InVision Studio:} Công cụ chuyên sâu về animation và chuyển động tương tác.
\item \textbf{Adobe Photoshop \& Illustrator:} Công cụ xử lý hình ảnh bitmap và đồ họa vector nền tảng.
\item \textbf{Affinity Designer, GIMP, Blender:} Các công cụ thiết kế bổ trợ, từ chỉnh sửa đồ họa 2D đến mô phỏng 3D hiện đại.
\end{itemize}

\begin{figure}[htbp]
\centering
\includegraphics[width=0.75\textwidth]{images/image3.png}
\caption{Các công cụ thiết kế giao diện đồ họa phổ biến}
\end{figure}
\begin{figure}[htbp]
\centering
\includegraphics[width=0.75\textwidth]{images/image4.png}
\caption{Hệ sinh thái công cụ UI/UX chuyên nghiệp}
\end{figure}

\bigskip\hrule\bigskip

\section{1.6 Bản chất khái niệm UI và UX}

Bốn trụ cột cơ bản của ngành thiết kế sản phẩm số tương tác:

\begin{itemize}[leftmargin=1.5em]
\item \textbf{Thiết kế tương tác (Interaction Design - IxD):} Tập trung vào cách thức con người tương tác qua lại với sản phẩm.
\item \textbf{Giao diện người dùng (User Interface - UI):} Điểm tiếp xúc trực quan giữa người và máy.
\item \textbf{Tương tác người - máy (Human-Computer Interaction - HCI):} Ngành khoa học nghiên cứu mối quan hệ tương hỗ giữa người dùng và máy tính.
\item \textbf{Kiến trúc thông tin (Information Architecture - IA):} Nghệ thuật tổ chức, sắp xếp và phân cấp cấu trúc thông tin của sản phẩm logic, trực quan.
\end{itemize}

\begin{verbatim}
       ┌────────────────────────────────────────────────────────┐
       │             TRẢI NGHIỆM NGƯỜI DÙNG (UX)                │
       │                                                        │
       │   ┌────────────────────────────────────────────────┐   │
       │   │           GIAO DIỆN NGƯỜI DÙNG (UI)            │   │
       │   │                                                │   │
       │   │  • Thiết kế trực quan (Visual Design)          │   │
       │   │    - Màu sắc, Typography, Icon, Layout         │   │
       │   │                                                │   │
       │   │  • Thiết kế tương tác (Interaction Design)     │   │
       │   │    - Hiệu ứng phản hồi, Chuyển động, Nút bấm   │   │
       │   └────────────────────────────────────────────────┘   │
       │                                                        │
       │  • Nghiên cứu người dùng (User Research)               │
       │  • Kiến trúc thông tin (Information Architecture)      │
       │  • Chiến lược nội dung & Kịch bản tương tác (Scenarios)│
       │  • Kiểm thử tính khả dụng (Usability Testing)          │
       └────────────────────────────────────────────────────────┘
\end{verbatim}

\bigskip\hrule\bigskip

\subsection{1.6.1 Giao diện người dùng (UI - User Interface)}
\begin{itemize}[leftmargin=1.5em]
\item \textbf{Bản chất của UI:} UI bao gồm hai thành phần mật thiết:
\end{itemize}
\begin{enumerate}[leftmargin=1.5em]
\item \textit{Thiết kế trực quan (Visual Design):} Hình dáng, màu sắc, font chữ, kích thước, khoảng cách, icon, bố cục màn hình.
\item \textit{Thiết kế tương tác (Interaction Design):} Cách thức hệ thống phản hồi lại các thao tác nhấp chuột, chạm, vuốt, cuộn của người dùng.
\end{enumerate}
\begin{itemize}[leftmargin=1.5em]
\item \textbf{Khái niệm:} Giao diện người dùng là không gian, là điểm truy cập nơi mà con người tiếp xúc và tương tác trực tiếp với một hệ thống hoặc sản phẩm.
\item \textbf{UI không chỉ giới hạn trong màn hình máy tính hay điện thoại:}
\item \textit{Ứng dụng/Website:} Màn hình, icon, thanh menu, nút click.
\item \textit{Trình duyệt Chrome:} Thanh URL, hệ thống các tab duyệt web, nút back/forward.
\item \textit{Máy bán nước tự động:} Hàng phím bấm chọn món, khe nhét tiền, hộc nhận lon nước.
\item \textit{Con dao:} Chiếc cán dao chính là giao diện để tay người cầm nắm và điều khiển lưỡi dao.
\item \textit{Băng ghế đá:} Bề mặt đá mà ta ngồi lên tiếp xúc chính là giao diện.
\item \textbf{Mục tiêu của UI Designer:} Tạo ra một giao diện thu hút, thẩm mỹ cao, trực quan và dễ dàng thao tác từ cấu trúc nội dung và chức năng của sản phẩm.
\end{itemize}

\subsubsection{Ví dụ về thiết kế UI tệ và UI tốt:}
\begin{itemize}[leftmargin=1.5em]
\item \textbf{Thiết kế UI tệ:}
\item Biểu mẫu nhập liệu dài dòng, cấu trúc rối rắm khiến khách hàng bỏ dở việc đăng ký.
\item Cỡ chữ quá bé, độ tương phản kém khiến mắt người đọc bị căng thẳng và khó tiếp thu thông tin.
\item Lạm dụng quá nhiều biểu tượng (icon) kỳ lạ, thiếu quy chuẩn gây hoang mang, bối rối cho người dùng.
\item \textbf{Thiết kế UI tốt:}
\item Ứng dụng đặt vé xem phim với quy trình 3 bước trực quan: Chọn phim $
\end{itemize}
ightarrow$ Chọn ghế $

ightarrow$ Thanh toán, hoàn tất trong chưa đầy 1 phút.

\begin{itemize}[leftmargin=1.5em]
\item Website thương mại điện tử hiện đại, phân cấp thông tin rõ ràng với nút Kêu gọi hành động (\textbf{CTA - Call to Action}) nổi bật, giúp khách hàng đưa ra quyết định nhanh chóng.
\end{itemize}

\bigskip\hrule\bigskip

\subsection{1.6.2 Trải nghiệm người dùng (UX - User Experience)}
\begin{itemize}[leftmargin=1.5em]
\item \textbf{Khái niệm:} Trải nghiệm người dùng là toàn bộ những nhận thức, cảm xúc, hành vi và phản ứng tâm lý mà một người trải qua trước, trong và sau quá trình tương tác với một sản phẩm hoặc dịch vụ.
\item \textbf{Tiêu chuẩn quốc tế ISO 9241-210 định nghĩa:} \textit{"User experience encompasses all aspects of the user's interaction with the company, its services, and its products. It is a person's perceptions and responses resulting from the use and/or anticipated use of a product, system or service."}
\item \textbf{Mục tiêu tối thượng của UX Design:}  
\end{itemize}
\[ \mathbf{ĐÁP\;ỨNG\;ĐƯỢC\;NHU\;CẦU\;CỦA\;NGƯỜI\;DÙNG} \]
\begin{itemize}[leftmargin=1.5em]
\item \textbf{UX trong đời sống hàng ngày:}
\item Khi cầm một chiếc kéo, việc cắt phải diễn ra trơn tru, sắc bén và không gây đau tay.
\item Khi ngồi lên một chiếc ghế đá công viên, cảm giác phải vững chãi, thoải mái, độ cao vừa vặn.
\item Khi dùng một chiếc smartphone, mọi thao tác vuốt chạm phải mượt mà, mang lại sự thỏa mãn, đáp ứng tức thì nhu cầu liên lạc và giải trí.
\item UX dùng công nghệ để thu hẹp khoảng cách giữa máy móc và con người, làm cho công nghệ trở nên thân thiện, nhân bản hơn.
\end{itemize}

\begin{figure}[htbp]
\centering
\includegraphics[width=0.75\textwidth]{images/image5.png}
\caption{Mục tiêu tối thượng của UX: Đáp ứng đúng nhu cầu người dùng}
\end{figure}
\begin{figure}[htbp]
\centering
\includegraphics[width=0.75\textwidth]{images/image6.png}
\caption{Trải nghiệm người dùng trong đời sống hàng ngày}
\end{figure}

\bigskip\hrule\bigskip

\subsection{1.6.3 Mối quan hệ tương hỗ giữa UI và UX}
\begin{itemize}[leftmargin=1.5em]
\item \textbf{Phân công trách nhiệm:}
\item UX Designer chịu trách nhiệm nghiên cứu hành vi, thấu hiểu tâm lý và phác họa nên toàn bộ \textbf{Hành trình khách hàng (Customer Journey)}.
\item UI Designer tập trung vào thiết kế từng chi tiết trực quan, từng điểm chạm cụ thể để hiện thực hóa hành trình đó một cách sống động, hấp dẫn.
\item \textbf{Nguyên lý quan trọng:} \textbf{UI LÀ MỘT PHẦN CỦA UX (UI is a component of UX).} Bạn không thể tạo ra một trải nghiệm người dùng trọn vẹn nếu thiếu đi một giao diện trực quan, rõ ràng.
\end{itemize}

\begin{figure}[htbp]
\centering
\includegraphics[width=0.75\textwidth]{images/image7.png}
\caption{Mối quan hệ hữu cơ giữa UI và UX}
\end{figure}
\begin{figure}[htbp]
\centering
\includegraphics[width=0.75\textwidth]{images/image8.png}
\caption{UI là một phần không thể tách rời trong tổng thể UX}
\end{figure}

\subsubsection{Tầm ảnh hưởng quyết định của UI/UX đối với sản phẩm số:}
\begin{itemize}[leftmargin=1.5em]
\item \textbf{Quy tắc 8 giây:} Các nghiên cứu chỉ ra rằng một sản phẩm số chỉ có khoảng \textbf{8 giây} đầu tiên để gây ấn tượng với khách hàng khi họ mở ứng dụng hoặc trang web.
\item \textbf{Quy tắc 75\% tin cậy:} Hơn \textbf{75\%} người tiêu dùng thừa nhận họ đánh giá mức độ uy tín và sự chuyên nghiệp của một doanh nghiệp dựa hoàn toàn vào tính thẩm mỹ và sự mượt mà của giao diện web/app.
\end{itemize}

\bigskip\hrule\bigskip

\subsection{1.6.4 Phân biệt CX, UX và BX}

Để có cái nhìn toàn diện trong môi trường doanh nghiệp, cần phân biệt 3 khái niệm:

\begin{itemize}[leftmargin=1.5em]
\item \textbf{UX (User Experience - Trải nghiệm người dùng):} Tập trung vào tương tác trực tiếp của người dùng với một sản phẩm số cụ thể (tính dễ dùng, tốc độ, luồng thao tác trên app/web).
\item \textbf{CX (Customer Experience - Trải nghiệm khách hàng):} Phạm vi rộng hơn, bao gồm tất cả các điểm chạm giữa khách hàng với doanh nghiệp xuyên suốt vòng đời khách hàng (từ quảng cáo, tư vấn, mua hàng, giao hàng, đổi trả đến dịch vụ hỗ trợ khách hàng).
\item \textbf{BX (Brand Experience - Trải nghiệm thương hiệu):} Cảm nhận tổng thể, giá trị cảm xúc và niềm tin lâu dài mà khách hàng dành cho toàn bộ thương hiệu.
\end{itemize}

\subsubsection{Case Study so sánh thực tế:}
\begin{enumerate}[leftmargin=1.5em]
\item \textbf{Trường hợp UX tốt nhưng CX kém:}  
\end{enumerate}
Khách hàng sử dụng ứng dụng đặt vé máy bay rất mượt mà, chọn chuyến và thanh toán trong 30 giây (UX tuyệt vời). Tuy nhiên, khi chuyến bay bị hoãn, khách hàng gọi lên tổng đài thì không ai bắt máy, ra quầy đổi vé thì thủ tục hành chính phức tạp, nhân viên cáu gắt (CX cực kỳ tồi tệ). Kết quả: Khách hàng vẫn bỏ sang hãng khác.

\begin{enumerate}[leftmargin=1.5em]
\item \textbf{Trường hợp UX chưa hoàn hảo nhưng CX và BX rất mạnh - Điển hình ZARA:}  
\end{enumerate}
Ứng dụng và website bán hàng của ZARA thường xuyên nhận phản hồi phàn nàn từ giới chuyên môn UX: màu sắc chữ khó đọc, ảnh người mẫu tạo dáng kỳ quặc, thanh điều hướng khác thường. Tuy nhiên, Zara lại có Trải nghiệm thương hiệu thời trang "ăn liền" (Fast Fashion) quá mạnh: mẫu mã ra mắt cực nhanh, giá cả hợp lý, hệ thống cửa hàng vật lý sang trọng ở vị trí đắc địa $

ightarrow$ Khách hàng vẫn mua sắm ồ ạt.

\begin{enumerate}[leftmargin=1.5em]
\item \textbf{Cách mạng hóa trải nghiệm gọi xe - Bài học từ UBER:}  
\end{enumerate}
Trước khi Uber xuất hiện, hành khách chịu đựng trải nghiệm khách hàng (CX) rất tệ: đứng vẫy taxi dưới trời mưa, tài xế gian lận đồng hồ tính tiền, thái độ phục vụ kém. Uber ra đời đã thay đổi toàn bộ cuộc chơi bằng việc cung cấp giải pháp UI/UX đột phá: biết trước giá tiền, theo dõi vị trí xe di chuyển thời gian thực trên bản đồ, thanh toán tự động, đánh giá chất lượng tài xế $

ightarrow$ Đưa trải nghiệm khách hàng lên một tầm cao mới.


\begin{figure}[htbp]
\centering
\includegraphics[width=0.75\textwidth]{images/image9.png}
\caption{ZARA: Thách thức trong thiết kế UX và sức mạnh của trải nghiệm thương hiệu}
\end{figure}

\bigskip\hrule\bigskip

\subsection{1.6.5 Các kịch bản tương tác giữa UI và UX trong thực tế}

Sản phẩm trong thực tế thường rơi vào các trạng thái sau:

\begin{itemize}[leftmargin=1.5em]
\item \textbf{Sản phẩm có UI tốt nhưng UX tệ:} Giao diện bóng bẩy, đồ họa 3D bắt mắt, màu sắc thời thượng nhưng cực kỳ khó tìm thông tin, nút bấm giấu kín, thao tác rườm rà. Người dùng mở ra khen đẹp nhưng dùng 2 phút thì phát bực và xóa app.
\item \textbf{Sản phẩm có UX tốt nhưng UI tệ:} Bố cục và tính năng cực kỳ tiện lợi, luồng đi chuẩn xác, nhưng font chữ xấu xí, màu sắc u ám, icon thô kệch. Người dùng vẫn dùng vì giải quyết được việc, nhưng cảm giác thiếu chuyên nghiệp.
\item \textbf{Sản phẩm đạt chuẩn hoàn hảo (UI tốt + UX tốt):} Vừa đẹp mắt vừa trực quan, dễ dùng, tạo ra sự thích thú và gắn bó lâu dài.
\end{itemize}

\subsubsection{Hai bài học thất bại đắt giá:}
\begin{itemize}[leftmargin=1.5em]
\item \textbf{Trường hợp 1 (UI phá hỏng ý tưởng kinh doanh xuất sắc):}  
\end{itemize}
Bạn xây dựng một ứng dụng với ý tưởng vô cùng sáng tạo, chưa ai làm. Nhưng khi phát hành, người dùng phàn nàn: font chữ quá bé, nút bấm quá nhỏ nằm sát nhau khiến họ liên tục bấm nhầm, hình ảnh vỡ hạt. Một thiết kế UI cẩu thả đã giết chết một ý tưởng kinh doanh triệu đô.

\begin{itemize}[leftmargin=1.5em]
\item \textbf{Trường hợp 2 (UX phức tạp phá hủy thương hiệu dịch vụ):}  
\end{itemize}
Một công ty đầu tư website với giao diện đồ họa cực kỳ hiện đại, bắt mắt. Tuy nhiên, khách hàng muốn đăng ký dịch vụ lại phải trải qua 7 trang biểu mẫu, không tìm thấy thông tin liên hệ hay mức giá. Khách hàng nản lòng rời đi sang đối thủ. Quy trình UX tồi đã trực tiếp gây thất thoát doanh thu nghiêm trọng.


\bigskip\hrule\bigskip

\section{1.7 Bảy quy tắc cơ bản đảm bảo tính UI/UX trong thiết kế}

\begin{enumerate}[leftmargin=1.5em]
\item \textbf{Thiết kế Đơn giản và Trực quan (Simplicity \& Intuitiveness):}  
\end{enumerate}
Giao diện cần đơn giản, dễ hiểu ngay từ cái nhìn đầu tiên. Đừng ép người dùng phải suy nghĩ hay học cách sử dụng một giao diện cơ bản.

\begin{enumerate}[leftmargin=1.5em]
\item \textbf{Tạo ra Trải nghiệm Người dùng Tích cực (Positive Experience):}  
\end{enumerate}
Đảm bảo người dùng đạt được mục đích của họ một cách nhanh chóng, thoải mái và không gặp phải các rào cản tâm lý.

\begin{enumerate}[leftmargin=1.5em]
\item \textbf{Đảm bảo Khả năng Tương tác (Interactivity):}  
\end{enumerate}
Các thành phần tương tác (nút bấm, thanh trượt, menu) phải phản hồi rõ ràng, mượt mà khi người dùng tác động.

\begin{enumerate}[leftmargin=1.5em]
\item \textbf{Đảm bảo Tính Tương thích (Responsiveness \& Cross-platform):}  
\end{enumerate}
Giao diện phải hiển thị tối ưu trên nhiều loại thiết bị, từ màn hình máy tính lớn, laptop đến máy tính bảng và điện thoại di động với các kích thước và tỷ lệ khung hình khác nhau.

\begin{enumerate}[leftmargin=1.5em]
\item \textbf{Tập trung vào Nội dung Chính (Content-First):}  
\end{enumerate}
Mọi chi tiết trang trí phải phục vụ cho việc làm nổi bật nội dung cốt lõi và luồng tác vụ chính của người dùng.

\begin{enumerate}[leftmargin=1.5em]
\item \textbf{Sử dụng Màu sắc và Hình ảnh Hợp lý (Color \& Imagery):}  
\end{enumerate}
Màu sắc phải tuân thủ hệ thống phân cấp thị giác (Visual Hierarchy), hỗ trợ khả năng đọc và phù hợp với ngữ cảnh văn hóa cũng như thương hiệu.

\begin{enumerate}[leftmargin=1.5em]
\item \textbf{Tạo Phong cách Đồ họa Thống nhất (Consistency in Graphic Style):}  
\end{enumerate}
Toàn bộ hệ thống phải đồng bộ về font chữ, kích cỡ nút, hệ thống icon, khoảng cách lề nhằm tạo nên một khối nhận diện hoàn chỉnh.


\bigskip\hrule\bigskip

\section{1.8 Phân loại giao diện người dùng theo hình thức tương tác}

\subsection{1.8.1 Giao diện WIMP (Windows, Icons, Menus, Pointers)}
Giao diện WIMP là mô hình tương tác đồ họa phổ biến nhất trong lịch sử điện toán cá nhân, được khởi xướng từ phòng thí nghiệm Xerox PARC, thương mại hóa thành công bởi Apple Macintosh và thống trị trên hệ điều hành Microsoft Windows:

\begin{itemize}[leftmargin=1.5em]
\item \textbf{Windows (Cửa sổ):} Các khu vực độc lập trên màn hình, hoạt động như các màn hình con. Người dùng có thể di chuyển, thay đổi kích thước, thu nhỏ, phóng to và sắp xếp nhiều cửa sổ làm việc cùng lúc. Cửa sổ đi kèm thanh tiêu đề (Title Bar), thanh cuộn (Scroll Bar) và các nút điều khiển thu phóng/đóng.
\item \textbf{Icons (Biểu tượng):} Các hình ảnh đồ họa trực quan nhỏ đại diện cho tệp tin, thư mục, chương trình hoặc thao tác. Giúp người dùng nhận diện nhanh chóng mà không cần đọc dòng lệnh.
\item \textbf{Menus (Trình đơn):} Danh sách các câu lệnh hoặc chức năng được hiển thị khi người dùng yêu cầu (Menu ngang, Menu thả xuống Dropdown, Menu ngữ cảnh Context Menu). Tên trên menu phải mang nghĩa rõ ràng.
\item \textbf{Pointers (Con trỏ):} Thiết bị trỏ (chuột, bàn di chuột touchpad, trackball, joystick) biểu thị vị trí thao tác trên màn hình, hình dạng con trỏ thay đổi theo ngữ cảnh (mũi tên, bàn tay click, con trỏ chữ I...).
\end{itemize}

\subsubsection{Các thành phần điều khiển gắn liền với WIMP:}
\begin{itemize}[leftmargin=1.5em]
\item \textbf{Radio Buttons:} Tập hợp các nút tròn cho phép người dùng chỉ chọn duy nhất 1 phương án trong một nhóm loại trừ lẫn nhau.
\item \textbf{Checkboxes:} Hộp kiểm cho phép người dùng chọn một, nhiều hoặc không chọn phương án nào.
\item \textbf{Toolbars:} Thanh công cụ chứa các biểu tượng thao tác nhanh thường dùng (Save, Undo, Bold...).
\item \textbf{Dialogue Boxes (Hộp thoại):} Cửa sổ nhỏ bật lên để yêu cầu người dùng xác nhận thông tin quan trọng hoặc cảnh báo lỗi trước khi tiếp tục.
\item \textbf{Palettes:} Bảng chọn công cụ hoặc bảng màu nổi trên vùng làm việc.
\end{itemize}

\bigskip\hrule\bigskip

\subsection{1.8.2 Thao tác trực tiếp (Direct Manipulation)}
\begin{itemize}[leftmargin=1.5em]
\item Người dùng tương tác trực tiếp lên các đối tượng ảo trên màn hình bằng cách kéo, thả, xoay, thay đổi kích thước, phóng to thu nhỏ tương tự như thao tác với đồ vật ngoài đời thực.
\item Thao tác trực tiếp mang lại phản hồi ngay lập tức, giảm thiểu tối đa gánh nặng ghi nhớ và rất tự nhiên đối với người sử dụng.
\end{itemize}

\begin{figure}[htbp]
\centering
\includegraphics[width=0.75\textwidth]{images/image10.png}
\caption{Thao tác trực tiếp (Direct Manipulation)}
\end{figure}

\bigskip\hrule\bigskip

\subsection{1.8.3 Giao diện Menu (Menu-driven Interface)}
\begin{itemize}[leftmargin=1.5em]
\item Cung cấp một tập hợp các lựa chọn có sẵn trên màn hình, người dùng lựa chọn bằng chuột, phím số hoặc phím điều hướng.
\item \textbf{Ưu điểm:}
\item Người dùng không phải ghi nhớ câu lệnh cú pháp vì mọi lựa chọn đều hiện diện trước mắt (Gợi ý thay vì nhớ lại).
\item Giới hạn không gian lựa chọn, tránh được việc người dùng nhập dữ liệu sai hoặc ngoại lệ.
\item Rút ngắn thời gian đào tạo cho người dùng mới.
\item \textbf{Nhược điểm:}
\item Chiếm diện tích không gian hiển thị trên màn hình.
\item Có thể gây chậm chạp và ức chế đối với người dùng chuyên nghiệp khi phải duyệt qua nhiều tầng menu con để thực hiện một tác vụ đơn giản.
\end{itemize}

\begin{figure}[htbp]
\centering
\includegraphics[width=0.75\textwidth]{images/image11.png}
\caption{Giao diện Menu (Menu-driven)}
\end{figure}

\bigskip\hrule\bigskip

\subsection{1.8.4 Giao diện dòng lệnh (Command Line Interface - CLI)}
\begin{itemize}[leftmargin=1.5em]
\item Người dùng tương tác với hệ điều hành thông qua việc gõ các câu lệnh bằng văn bản chính xác theo đúng cú pháp (ví dụ: màn hình MS-DOS, Command Prompt, PowerShell, Linux Terminal).
\item Một số lệnh cơ bản: \texttt{dir} (xem nội dung thư mục trên DOS), \texttt{ls} (trên Linux), \texttt{cd} (chuyển thư mục)...
\item \textbf{Ưu điểm:}
\item Cực kỳ linh hoạt, tốc độ thực thi rất nhanh khi người dùng đã thành thạo.
\item Tiêu tốn cực kỳ ít tài nguyên phần cứng (RAM, CPU, đồ họa).
\item Dễ dàng lập trình tự động hóa các tác vụ phức tạp thông qua kịch bản (Scripting).
\item \textbf{Nhược điểm:}
\item Rất khó học đối với người dùng phổ thông.
\item Đòi hỏi người dùng phải ghi nhớ chính xác câu lệnh, tùy chọn tham số và cú pháp khắt khe.
\item Phản hồi trực quan hạn chế, rất dễ gây ra sai sót nghiêm trọng nếu gõ nhầm lệnh.
\end{itemize}

\begin{figure}[htbp]
\centering
\includegraphics[width=0.75\textwidth]{images/image12.png}
\caption{Giao diện dòng lệnh (Command Line Interface - CLI)}
\end{figure}

\bigskip\hrule\bigskip

\subsection{1.8.5 Giao diện dựa trên biểu mẫu (Form-based / Fill-in forms)}
\begin{itemize}[leftmargin=1.5em]
\item Giao diện cung cấp các trường dữ liệu được thiết kế sẵn (Text box, Date picker, Dropdown...) để người dùng điền thông tin thích hợp vào.
\item Được sử dụng phổ biến nhất trong các hệ thống thông tin quản lý, đăng ký tài khoản, giao dịch ngân hàng và cài đặt phần mềm.
\item Đòi hỏi giao diện phải có cơ chế xác thực dữ liệu đầu vào (Validation) tức thì và thông báo lỗi rõ ràng.
\end{itemize}

\bigskip\hrule\bigskip

\subsection{1.8.6 Giao diện ngôn ngữ tự nhiên (Natural Language Interface)}
\begin{itemize}[leftmargin=1.5em]
\item Cho phép người dùng giao tiếp với hệ thống bằng chính ngôn ngữ nói hoặc viết hàng ngày của con người (ví dụ: Chatbot AI, Trợ lý ảo ChatGPT, Gemini, Siri, Google Assistant).
\item \textbf{Ưu điểm:} Vô cùng tự nhiên, không đòi hỏi người dùng phải học bất kỳ cú pháp kỹ thuật nào.
\item \textbf{Thách thức:} Ngôn ngữ tự nhiên có tính đa nghĩa, phụ thuộc vào ngữ cảnh và vùng miền, đòi hỏi hệ thống AI phải xử lý ngôn ngữ tự nhiên (NLP) cực kỳ phức tạp.
\end{itemize}

\bigskip\hrule\bigskip

\subsection{1.8.7 Giao diện ngôn ngữ truy vấn (Query Language - SQL)}
\begin{itemize}[leftmargin=1.5em]
\item Được sử dụng chuyên biệt để khai phá, truy vấn và quản lý dữ liệu trong hệ thống cơ sở dữ liệu quan hệ (RDBMS).
\item Đòi hỏi người sử dụng phải hiểu sâu về cấu trúc cơ sở dữ liệu (Database Schema, Tables, Relations) và cú pháp câu lệnh (SQL: \texttt{SELECT}, \texttt{INSERT}, \texttt{UPDATE}, \texttt{WHERE}, \texttt{JOIN}...).
\item Phù hợp cho lập trình viên, kỹ sư dữ liệu và chuyên viên phân tích hệ thống.
\end{itemize}

\bigskip\hrule\bigskip

\subsection{1.8.8 Các phong cách tương tác tiên tiến: VR, AR, Haptic và Cử chỉ}
\begin{itemize}[leftmargin=1.5em]
\item \textbf{Thực tế ảo (Virtual Reality - VR):} Đưa người dùng đắm chìm hoàn toàn vào một không gian mô phỏng 3D nhân tạo, tạo cảm giác như đang hiện diện trong môi trường thực.
\item \textbf{Thực tế tăng cường (Augmented Reality - AR):} Cầu nối giữa thế giới thực và kỹ thuật số bằng cách đặt chồng các vật thể ảo, thông tin đồ họa 3D lên hình ảnh thế giới vật lý thực thông qua camera điện thoại hoặc kính thông minh.
\item \textbf{Phản hồi xúc giác (Haptic Interface):} Cung cấp phản hồi cảm giác chạm bằng cách sử dụng độ rung, áp lực hoặc lực cản cơ học (ví dụ: tay cầm chơi game DualSense, bàn rê chuột Force Touch trên MacBook).
\item \textbf{Giao diện cử chỉ (Gesture-based Interface):} Cho phép người dùng điều khiển hệ thống thông qua các chuyển động tự nhiên của bàn tay, ngón tay hoặc cơ thể (ví dụ: vẫy tay chuyển bài hát, cảm biến Kinect, Apple Vision Pro).
\end{itemize}

\begin{figure}[htbp]
\centering
\includegraphics[width=0.75\textwidth]{images/image13.png}
\caption{Các phong cách tương tác hiện đại: VR, AR, Haptic, Cử chỉ}
\end{figure}

\bigskip\hrule\bigskip

\section{1.9 Các nguyên lý thiết kế giao diện người dùng nền tảng}

Thiết kế giao diện người dùng phải đặt trọng tâm vào nhu cầu, kinh nghiệm và năng lực tiếp nhận của người dùng. Chín nguyên lý nền tảng bao gồm:

\begin{enumerate}[leftmargin=1.5em]
\item \textbf{Quen thuộc với người dùng (User Familiarity):} Sử dụng các khái niệm, biểu tượng và thuật ngữ bắt nguồn từ kinh nghiệm thực tế đời thường của người dùng thay vì dùng thuật ngữ kỹ thuật chuyên môn.
\item \textbf{Nhất quán (Consistency):} Giao diện phải duy trì sự đồng bộ từ cách bố trí, phím tắt, biểu tượng cho đến luồng thao tác trên toàn bộ hệ sinh thái sản phẩm.
\item \textbf{Tối thiểu hóa sự bất ngờ (Minimal Surprise):} Hành vi của hệ thống phải diễn ra đúng như mong đợi tự nhiên của người dùng, không bao giờ gây hoang mang, khó hiểu.
\item \textbf{Khả năng khôi phục (Recoverability):} Hệ thống phải cung cấp cho người dùng con đường cứu vãn khi họ thao tác nhầm lẫn (hỗ trợ Undo, Redo, hoặc yêu cầu xác nhận trước khi thực hiện hành động phá hủy dữ liệu).
\item \textbf{Hướng dẫn người dùng (User Guidance):} Cung cấp phản hồi kịp thời, thông báo lỗi mang tính xây dựng và các gợi ý hướng dẫn đúng thời điểm theo ngữ cảnh.
\item \textbf{Tính đa dạng người dùng (User Diversity):} Thiết kế phải tính đến nhiều đối tượng người dùng khác nhau (người mới làm quen dùng menu đồ họa; người chuyên nghiệp sử dụng phím tắt nâng cao).
\item \textbf{Thiết kế trực quan (Visual Clarity):} Giao diện có bố cục rõ ràng, nhìn vào là hiểu ngay cách dùng trong lần đầu tiên tiếp xúc mà không cần phải trải qua đào tạo phức tạp.
\item \textbf{Quy tắc 3-Clicks (The 3-Click Rule):} Người dùng nên tìm thấy thông tin hoặc thực hiện được chức năng mong muốn chỉ trong vòng tối đa \textbf{3 lần nhấp chuột}.
\item \textbf{Hiện tượng tâm lý học trong thiết kế:}
\end{enumerate}
\begin{itemize}[leftmargin=1.5em]
\item \textit{Hội chứng "gà con" (Chick Imprinting Syndrome):} Con người có xu hướng yêu thích và gắn bó với hệ thống đầu tiên mà họ đã quen sử dụng, đồng thời có tâm lý e ngại hoặc phản kháng đối với những hệ thống có cách tương tác quá khác biệt.
\item \textit{Xu hướng thị giác:} Phần lớn người dùng có xu hướng thích nhìn biểu tượng trực quan hơn là phải đọc những dòng văn bản dài dòng.
\end{itemize}

\subsubsection{Ví dụ thực tế về tối ưu hóa cấu trúc Menu danh mục sản phẩm:}
\begin{itemize}[leftmargin=1.5em]
\item Khi tổ chức danh mục hàng hóa theo cây đa cấp 3 tầng truyền thống, người dùng phải mất ít nhất 2 lần di chuột và 3 lần tìm kiếm để chọn được mục mong muốn (ví dụ "Nồi cơm điện").
\item Bằng cách tối ưu hóa: Chuyển sang cây danh mục 2 cấp dàn ngang, tận dụng tối đa khoảng trống của màn hình phía bên phải để hiển thị toàn bộ các nhóm hàng hóa. Người dùng chỉ cần 1 lần rê chuột và tìm kiếm tức thì.
\end{itemize}

\begin{figure}[htbp]
\centering
\includegraphics[width=0.75\textwidth]{images/image14.png}
\caption{Cấu trúc menu đa cấp truyền thống}
\end{figure}
\begin{figure}[htbp]
\centering
\includegraphics[width=0.75\textwidth]{images/image15.png}
\caption{Cách phân vùng thông tin menu tối ưu}
\end{figure}
\begin{figure}[htbp]
\centering
\includegraphics[width=0.75\textwidth]{images/image16.png}
\caption{Menu dàn ngang tận dụng không gian màn hình}
\end{figure}

\bigskip\hrule\bigskip

\section{1.10 Phương pháp thiết kế giao diện theo 3 bước}

Quy trình thiết kế một giao diện chuẩn mực được thực hiện qua 3 bước bài bản:

\begin{enumerate}[leftmargin=1.5em]
\item \textbf{Bước 1: Xác định thông tin:}  
\end{enumerate}
\begin{itemize}[leftmargin=1.5em]
\item Phân tích sâu các mô hình dữ liệu (Mô hình thực thể kết hợp ERD, Lược đồ cơ sở dữ liệu quan hệ).  
\item Xác định chính xác các trường dữ liệu cần nhập, hiển thị hoặc tra cứu trên giao diện.
\end{itemize}
\begin{enumerate}[leftmargin=1.5em]
\item \textbf{Bước 2: Thiết kế giao diện đúng (Functional Correctness):}  
\end{enumerate}
\begin{itemize}[leftmargin=1.5em]
\item Đảm bảo giao diện có đầy đủ chức năng và đúng đắn về mặt logic nghiệp vụ.  
\item Ở bước này, mục tiêu số một là \textbf{đúng chức năng}, chưa cần tối ưu sâu về tính tiện dụng hay thẩm mỹ.
\end{itemize}
\begin{enumerate}[leftmargin=1.5em]
\item \textbf{Bước 3: Tinh chỉnh giao diện tiện dụng (Usability Refinement):}  
\end{enumerate}
\begin{itemize}[leftmargin=1.5em]
\item Dựa trên phân tích hoạt động và hành vi thực tế của người dùng để sắp xếp lại các thành phần.  
\item Tối ưu hóa phân cấp thị giác, rút ngắn quãng đường di chuyển chuột, hỗ trợ phím Tab và nâng cao tính khả dụng tối đa.
\end{itemize}

\begin{figure}[htbp]
\centering
\includegraphics[width=0.75\textwidth]{images/image17.png}
\caption{Bước 1: Phân tích mô hình dữ liệu ERD}
\end{figure}
\begin{figure}[htbp]
\centering
\includegraphics[width=0.75\textwidth]{images/image18.png}
\caption{Bước 2: Thiết kế giao diện đảm bảo tính đúng đắn}
\end{figure}
\begin{figure}[htbp]
\centering
\includegraphics[width=0.75\textwidth]{images/image19.png}
\caption{Bước 3: Tinh chỉnh giao diện tối ưu hóa tính tiện dụng}
\end{figure}

\bigskip\hrule\bigskip

\section{1.11 Quy trình công việc chi tiết của UI Designer và UX Designer}

Trong một dự án phát triển phần mềm, vai trò của UI và UX được phân định rõ ràng nhưng kết hợp nhịp nhàng:


\subsection{1.11.1 Trách nhiệm của UX Designer (Nghiên cứu \& Cấu trúc)}
\begin{enumerate}[leftmargin=1.5em]
\item Tiếp nhận và tìm hiểu sâu sắc mục tiêu chiến lược của dự án.
\item Xác định nhóm đối tượng người dùng trọng tâm của sản phẩm.
\item Thu thập yêu cầu và phân tích sâu các nhu cầu tiềm ẩn của người dùng.
\item Mô hình hóa quy trình thao tác nghiệp vụ thông qua các sơ đồ luồng (Flowcharts, User Flows).
\item Xây dựng các chân dung khách hàng (\textbf{Personas}), kịch bản sử dụng (\textbf{Scenarios}) và dự báo các rủi ro tương tác.
\item Xây dựng các bản phác thảo khung xương (\textbf{Mockup / Lo-fi Wireframe}) và tiến hành kiểm tra sớm với người dùng.
\item Xây dựng nguyên mẫu thử nghiệm (\textbf{Prototype}) để thử nghiệm thực tế.
\item Tổ chức kiểm thử chấp nhận người dùng (\textbf{User Acceptance Testing - UAT}) và phân tích kết quả.
\item Bàn giao cấu trúc và quy chuẩn cho đội ngũ UI và Developers.
\end{enumerate}

\subsection{1.11.2 Trách nhiệm của UI Designer (Trực quan \& Hiện thực hóa)}
\begin{enumerate}[leftmargin=1.5em]
\item Định hình ngôn ngữ thiết kế thị giác (Visual Design Language) của sản phẩm.
\item Lựa chọn bảng màu, phông chữ, biểu tượng, khoảng cách và phong cách hình ảnh nhất quán.
\item Đảm bảo tính đồng nhất trên toàn hệ thống: Ví dụ, nếu nút Đóng được chọn màu đỏ bo góc 8px thì tất cả các nút Đóng trong toàn bộ ứng dụng phải tuân theo đúng quy chuẩn đó, tuyệt đối không được tùy tiện thay đổi.
\item Thiết kế các trạng thái tương tác của thành phần: Normal, Hover, Focus, Active, Disabled.
\item Xây dựng hệ thống UI Kit và Design System hoàn chỉnh phục vụ cho việc mở rộng sản phẩm lâu dài.
\end{enumerate}

\begin{figure}[htbp]
\centering
\includegraphics[width=0.75\textwidth]{images/image20.png}
\caption{So sánh vai trò và công việc giữa UX Designer và UI Designer}
\end{figure}

\bigskip\hrule\bigskip

\section{1.12 Quy trình thiết kế giao diện lặp (Iterative Design Process)}

Thiết kế giao diện người dùng không phải là một đường thẳng mà là một \textbf{quy trình lặp đi lặp lại (Iterative Process)} bao gồm 3 hoạt động cốt lõi:

\begin{enumerate}[leftmargin=1.5em]
\item \textbf{Phân tích người sử dụng (User Analysis):} Thấu hiểu sâu sắc người dùng sẽ làm gì với hệ thống.
\item \textbf{Lập mẫu thử hệ thống (System Prototyping):} Xây dựng các phiên bản mô phỏng để người dùng trải nghiệm trực tiếp.
\item \textbf{Đánh giá giao diện (Interface Evaluation):} Thử nghiệm và thu thập phản hồi để phát hiện các lỗ hổng thiết kế và tiếp tục vòng lặp cải tiến.
\end{enumerate}

\begin{figure}[htbp]
\centering
\includegraphics[width=0.75\textwidth]{images/image21.png}
\caption{Vòng lặp quy trình thiết kế giao diện người dùng}
\end{figure}

\subsection{Các kỹ thuật phân tích người dùng:}
\begin{itemize}[leftmargin=1.5em]
\item \textbf{Phân tích nhiệm vụ (Task Analysis):} Mô hình hóa chi tiết từng bước mà người dùng phải thực hiện để hoàn thành một mục tiêu cụ thể.
\item \textbf{Phân tích nhiệm vụ phân cấp (Hierarchical Task Analysis - HTA):} Chia nhỏ một tác vụ lớn thành các nhiệm vụ con theo từng cấp độ logic.
\item \textbf{Phỏng vấn và Trắc nghiệm (Interviews \& Questionnaires):} Đặt câu hỏi trực tiếp hoặc qua biểu mẫu khảo sát để nắm bắt mong muốn và cảm xúc của người dùng.
\item \textbf{Mô tả và Quan sát (Observation):} Trực tiếp quan sát người dùng thao tác trong môi trường làm việc thực tế của họ để phát hiện những thói quen mà họ không tự nói ra được.
\end{itemize}

\subsection{Kỹ thuật lập mẫu thử (Prototyping):}
\begin{itemize}[leftmargin=1.5em]
\item \textbf{Giai đoạn 1 - Lập mẫu thử trên giấy (Paper Prototype):} Phác thảo nhanh bằng bút trên giấy để thử nghiệm luồng ý tưởng sơ khởi với chi phí cực rẻ và tốc độ nhanh.
\item \textbf{Giai đoạn 2 - Xây dựng và tinh chỉnh mẫu thử kỹ thuật số (Digital Prototype):} Sử dụng các công cụ chuyên dụng (Figma, Protopie, hoặc trước đây là Macromedia Director, Visual Basic, Web scripts) để mô phỏng tương tác như sản phẩm thật.
\end{itemize}

\subsection{Kỹ thuật đánh giá giao diện người dùng:}
\begin{itemize}[leftmargin=1.5em]
\item Khảo sát trắc nghiệm ý kiến người sử dụng.
\item Cài đặt mã lệnh theo dõi (Instrumentation / Analytics) trong phần mềm để tự động ghi lại hành vi click, thời gian dừng màn hình và tỷ lệ bỏ trang của người dùng trực tuyến.
\end{itemize}

\begin{figure}[htbp]
\centering
\includegraphics[width=0.75\textwidth]{images/image22.png}
\caption{Kỹ thuật đánh giá giao diện người dùng}
\end{figure}

\bigskip\hrule\bigskip

\section{1.13 Xây dựng chân dung người dùng (User Persona)}

\subsection{1.13.1 Persona là gì?}
\textbf{Persona (Chân dung người dùng giả định)} là một hồ sơ mô tả bán hư cấu về người dùng lý tưởng của sản phẩm, được xây dựng dựa trên dữ liệu nghiên cứu người dùng thực tế và các phân tích định lượng/định tính. Persona mang lại cảm giác chân thực như một con người bằng xương bằng thịt với tên tuổi, hình ảnh, nghề nghiệp, sở thích và mục tiêu rõ ràng.


\subsection{1.13.2 Ba loại Persona chính trong UX:}
\begin{enumerate}[leftmargin=1.5em]
\item \textbf{Proto-persona:} Chân dung khách hàng được tạo ra nhanh chóng dựa trên các giả định ban đầu và kinh nghiệm của nhóm nội bộ, chưa trải qua quá trình nghiên cứu thực địa sâu.
\item \textbf{Marketing Persona:} Chân dung tập trung vào mục đích mua sắm, kênh tiếp cận, hành vi tiêu dùng và động lực thúc đẩy khách hàng chi tiền.
\item \textbf{Design Persona:} Chân dung chuyên sâu phục vụ cho thiết kế UX/UI, tập trung vào mục tiêu thao tác, bối cảnh sử dụng sản phẩm và các rào cản (pain points) mà người dùng gặp phải khi sử dụng ứng dụng.
\end{enumerate}

\bigskip\hrule\bigskip

\subsection{1.13.3 Tại sao Persona lại quan trọng?}
Doanh nghiệp xây dựng sản phẩm là để giải quyết các vấn đề cụ thể cho các nhóm người dùng cụ thể. Một trong những sai lầm chết người nhất là thiết kế một sản phẩm "cho tất cả mọi người", bởi khi cố gắng làm hài lòng tất cả mọi người, sản phẩm sẽ chẳng phục vụ tốt cho bất kỳ ai.


\begin{myquote}
\textbf{Tầm quan trọng của Persona:}
- Xác định rõ ràng: \textit{Người dùng của bạn là ai?}
- Và quan trọng không kém: \textit{Người dùng của bạn KHÔNG PHẢI là ai?}
\end{myquote}

\bigskip\hrule\bigskip

\subsection{1.13.4 Bốn thành phần cốt lõi trong một bản User Persona}
Một form Persona tiêu chuẩn (như chuẩn thiết kế của Google) thường bao gồm:

\begin{enumerate}[leftmargin=1.5em]
\item \textbf{Thông tin nhân khẩu học (Demographics):} Tên đầy đủ, hình ảnh chân dung đại diện, tuổi tác, giới tính, tình trạng hôn nhân, trình độ học vấn, nghề nghiệp, mức thu nhập.
\item \textbf{Mục đích của Persona (Goals / Motivations):} Điều gì thôi thúc người này tìm đến và sử dụng sản phẩm? Họ muốn đạt được kết quả cụ thể nào?
\item \textbf{Cản trở của Persona (Frustrations / Pain Points):} Những khó khăn, rắc rối, ức chế nào đang cản trở họ đạt được mục tiêu khi sử dụng các giải pháp hiện tại?
\item \textbf{Mô tả / Câu trích dẫn (Quotes / Bio):} Những phát biểu điển hình phản ánh thái độ sống, tính cách và quan điểm của nhóm người dùng đại diện.
\end{enumerate}

\begin{figure}[htbp]
\centering
\includegraphics[width=0.75\textwidth]{images/image23.png}
\caption{Cấu trúc các thành phần trong một User Persona}
\end{figure}
\begin{figure}[htbp]
\centering
\includegraphics[width=0.75\textwidth]{images/image24.png}
\caption{Mẫu User Persona chuẩn từ Google}
\end{figure}

\bigskip\hrule\bigskip

\subsection{1.13.5 Cần phải xây dựng bao nhiêu Persona cho một sản phẩm?}
Không có một con số tuyệt đối, tuy nhiên thông thường một sản phẩm nên tập trung xây dựng từ \textbf{2 đến 3 Persona} chính. Quá nhiều Persona sẽ làm phân tán nguồn lực và gây mâu thuẫn trong việc ưu tiên tính năng.


\subsection{1.13.6 Lợi ích cốt lõi của việc sử dụng Persona}
\begin{enumerate}[leftmargin=1.5em]
\item \textbf{Công cụ giao tiếp và định hướng nội bộ hiệu quả:}  
\end{enumerate}
User Persona tạo nên một "ngôn ngữ chung" thống nhất giữa đội ngũ sản phẩm, lập trình viên, marketing và ban lãnh đạo. Khi thảo luận tính năng mới, nhóm sẽ không tranh cãi theo cảm tính cá nhân ("Tôi nghĩ tính năng này hay"), mà sẽ dựa trên Persona ("Liệu anh Tuấn 45 tuổi có hiểu và dùng tính năng này không?").

\begin{enumerate}[leftmargin=1.5em]
\item \textbf{Kim chỉ nam cho mọi quyết định của sản phẩm (User-Centric Compass):}  
\end{enumerate}
\begin{itemize}[leftmargin=1.5em]
\item Ở khâu lên ý tưởng: Persona giúp chọn lọc giải pháp phù hợp nhất.
\item Ở khâu kiểm thử (Testing): Persona là tiêu chuẩn để tuyển chọn đúng đối tượng tham gia thử nghiệm.
\item Ở khâu truyền thông: Persona giúp nhóm Marketing nhắm đúng đối tượng và thông điệp.
\end{itemize}

\begin{figure}[htbp]
\centering
\includegraphics[width=0.75\textwidth]{images/image25.png}
\caption{Tầm quan trọng và lợi ích của việc xác định User Persona}
\end{figure}

\bigskip\hrule\bigskip

\section{1.14 Sơ đồ trang web (Sitemap) và Luồng công việc (Workflow)}

\subsection{1.14.1 Sơ đồ trang web (Sitemap)}
\begin{itemize}[leftmargin=1.5em]
\item \textbf{Định nghĩa:} Sitemap là một sơ đồ hoặc danh sách liệt kê phân cấp tất cả các trang, nội dung, hình ảnh và tài nguyên của một website.
\item \textbf{Lợi ích:}
\item Giúp các công cụ tìm kiếm (Google, Bing) thu thập dữ liệu và lập chỉ mục (index) website nhanh chóng, chính xác.
\item Định rõ những URL quan trọng được ưu tiên hiển thị.
\item Tối ưu hóa thời gian và chi phí SEO, mang lại sự thuận tiện khi tìm kiếm cho người dùng.
\item \textbf{Phân loại theo cấu trúc:}
\item \textit{HTML Sitemap:} Dành cho người dùng điều hướng trực tiếp trên trang web.
\item \textit{XML Sitemap:} Dành riêng cho robot của các công cụ tìm kiếm đọc và lập chỉ mục.
\item \textbf{Phân loại theo định dạng nội dung:}
\item \textit{Image Sitemap:} Chứa thông tin các hình ảnh lưu trữ trên website.
\item \textit{Video Sitemap:} Tổng hợp các thông tin về video.
\item \textit{News Sitemap:} Phục vụ cho Google News thu thập thông tin báo chí nhanh chóng.
\item \textit{Mobile Sitemap:} Cần thiết khi website có các trang chuyên biệt hiển thị trên thiết bị di động.
\end{itemize}

\bigskip\hrule\bigskip

\subsection{1.14.2 Luồng công việc (Workflow)}
\begin{itemize}[leftmargin=1.5em]
\item \textbf{Định nghĩa:} Workflow là sơ đồ biểu diễn thứ tự tuần tự từng bước công việc, sự kiện hoặc quyết định cần thực hiện để hoàn thành một quy trình nghiệp vụ cụ thể.
\item \textbf{Ví dụ về một luồng mua - bán hàng hóa trực tuyến:}
\end{itemize}
\[ \text{Xem sản phẩm & giá} \longrightarrow \text{Chọn mua} \longrightarrow \text{Điền form thông tin} \longrightarrow \text{Thanh toán đơn hàng} \]
\begin{itemize}[leftmargin=1.5em]
\item \textbf{Các phần mềm thiết kế Workflow phổ biến:} Microsoft Project, Microsoft Excel, Microsoft Visio, Workflow Designer trong Visual Studio.
\item \textbf{Lý do nên sử dụng Workflow:}
\item Mô hình hóa trực quan và đơn giản hóa các luồng điều khiển nghiệp vụ phức tạp.
\item Cho phép kiểm tra tính đúng đắn của logic phần mềm, theo dõi vết thực thi thời gian thực.
\item Tạo cầu nối phối hợp nhịp nhàng giữa con người và hệ thống phần mềm.
\end{itemize}

\begin{figure}[htbp]
\centering
\includegraphics[width=0.75\textwidth]{images/image26.png}
\caption{Quy trình luồng công việc (Workflow)}
\end{figure}
\begin{figure}[htbp]
\centering
\includegraphics[width=0.75\textwidth]{images/image27.png}
\caption{Sơ đồ chi tiết quy trình xử lý công việc}
\end{figure}
\begin{figure}[htbp]
\centering
\includegraphics[width=0.75\textwidth]{images/image28.png}
\caption{Mô hình hóa các bước thao tác trong hệ thống}
\end{figure}

\bigskip\hrule\bigskip

\section{1.15 Bố cục website tiêu chuẩn và Case Study "Nút bấm 300 triệu USD"}

\subsection{1.15.1 Bố cục tiêu chuẩn của một trang web (Standard Web Layout)}
Một bố cục website kinh điển gồm 5 khu vực chức năng:

\begin{enumerate}[leftmargin=1.5em]
\item \textbf{Header (Đầu trang):} Phần đầu tiên của trang web, chứa Logo thương hiệu, banner chính, thanh tìm kiếm và menu phụ (đăng nhập, ngôn ngữ, liên hệ).
\item \textbf{Navigation (Thanh điều hướng / Menu chính):} Thanh ngang chứa các thể loại hoặc chuyên mục chính của website.
\item \textbf{Content (Nội dung chính):} Khu vực chiếm diện tích lớn nhất, hiển thị nội dung trọng tâm của trang hiện tại.
\item \textbf{Sidebar (Thanh bên):} Nằm bên trái hoặc bên phải nội dung chính, chứa các menu phụ, bài viết liên quan, bộ lọc sản phẩm hoặc quảng cáo.
\item \textbf{Footer (Chân trang):} Nằm ở đáy trang, chứa thông tin bản quyền (Copyright), chính sách bảo mật, điều khoản sử dụng, thông tin công ty và liên kết mạng xã hội.
\end{enumerate}

\begin{figure}[htbp]
\centering
\includegraphics[width=0.75\textwidth]{images/image29.png}
\caption{Bố cục website tiêu chuẩn (Header, Nav, Content, Sidebar, Footer)}
\end{figure}
\begin{figure}[htbp]
\centering
\includegraphics[width=0.75\textwidth]{images/image30.png}
\caption{Các biến thể bố cục website hiện đại}
\end{figure}
\begin{figure}[htbp]
\centering
\includegraphics[width=0.75\textwidth]{images/image31.png}
\caption{Hệ thống lưới trong phân chia bố cục giao diện}
\end{figure}

\bigskip\hrule\bigskip

\subsection{1.15.2 Case Study kinh điển: "Chiếc nút bấm trị giá 300 triệu USD"}
\textit{(Theo cuốn sách kinh điển về thiết kế giao diện: "Don't Make Me Think" - Steve Krug và chuyên gia UX Jared Spool)}


\subsubsection{Quy tắc vàng:}
\begin{myquote}
\textbf{"Người dùng rất lười, hãy thiết kế sao cho họ ít phải suy nghĩ nhất!"}
Thiết kế xuất sắc không chỉ nằm ở vẻ đẹp bóng bẩy, mà là sự đơn giản tối đa, dễ sử dụng nhất và giúp người dùng tốn ít công sức nhất.
\end{myquote}

\subsubsection{Câu chuyện thực tế:}
\begin{itemize}[leftmargin=1.5em]
\item \textbf{Bối cảnh:} Một trang web thương mại điện tử lớn có đầy đủ tính năng: hiển thị hàng hóa, cho hàng vào giỏ và thanh toán.
\item \textbf{Rào cản chết người:} Sau khi khách hàng đã chọn hàng đầy giỏ và nhấn "Thanh toán", một form đăng nhập nhỏ xuất hiện gồm 2 ô: Username, Password, cùng 2 nút: \textit{Login} (Đăng nhập) và \textit{Register} (Đăng ký), kèm link \textit{Quên mật khẩu}.
\item Đội ngũ Lập trình viên cho rằng đây là quy trình hoàn toàn hợp lý và mang lại sự tiện lợi cho những lần mua sắm sau.
\item Nhưng trong thực tế, chiếc form nhỏ này đã gây ra sự sụt giảm doanh thu thảm hại! Khách hàng mới không muốn mất thời gian điền form đăng ký dài dòng chỉ để mua một món đồ. Còn khách hàng cũ thì phần lớn đều quên mất mật khẩu của mình, việc nhận email khôi phục mật khẩu quá rườm rà khiến họ từ bỏ giỏ hàng ngay lập tức.
\item \textbf{Giải pháp tinh tế:}
\item Nhóm chuyên gia UX đã đề xuất một thay đổi vô cùng đơn giản: \textbf{Xóa bỏ hoàn toàn nút Register}.
\item Thay vào đó là nút \textbf{"Continue" (Tiếp tục)} kèm theo lời giải thích thân thiện: \textit{"Bạn không cần tạo tài khoản để mua sắm. Hãy bấm nút Tiếp tục để thanh toán ngay dưới dạng khách (Guest Checkout). Để việc thanh toán nhanh hơn ở những lần sau, bạn có thể tạo tài khoản tùy chọn sau khi đã hoàn tất đơn hàng."}
\item \textbf{Kết quả kinh ngạc:}
\item Số lượng khách hàng hoàn tất đơn hàng tăng vọt \textbf{45\%}.
\item Trong năm đầu tiên áp dụng, doanh số của trang web đã tăng thêm \textbf{300.000.000 USD} (300 triệu USD) chỉ nhờ thay đổi đúng một nút bấm!
\item Bài học: Thiết kế UI/UX không chỉ là việc trang trí bề ngoài, mà trực tiếp quyết định sự sống còn và doanh thu của cả một đế chế kinh doanh.
\end{itemize}


\bigskip\hrule\bigskip

\chapter{CHƯƠNG 2: CÁC NGUYÊN LÝ VÀ QUY TRÌNH THIẾT KẾ GIAO DIỆN}

\section{2.1 Tính khả dụng (Usability) và 5 thành phần cốt lõi}

\subsection{2.1.1 Định nghĩa về Usability}
\begin{itemize}[leftmargin=1.5em]
\item \textbf{Usability (Tính khả dụng / Khả năng sử dụng):} Là thước đo đánh giá mức độ dễ dàng, trơn tru và hiệu quả khi người dùng tương tác với một sản phẩm hoặc hệ thống.
\item Thiết kế đạt chuẩn Usability là thiết kế giúp người dùng hoàn thành mục tiêu của họ một cách tự nhiên nhất, tốn ít thời gian và công sức nhất.
\item \textit{"Usability is one of the most important concepts in UI and UX, and therefore interfaces as well."}
\end{itemize}

\begin{figure}[htbp]
\centering
\includegraphics[width=0.75\textwidth]{images/image32.png}
\caption{Khái niệm Usability trong thiết kế giao diện}
\end{figure}
\begin{figure}[htbp]
\centering
\includegraphics[width=0.75\textwidth]{images/image33.png}
\caption{Tầm quan trọng của Usability trong hệ sinh thái sản phẩm}
\end{figure}

\bigskip\hrule\bigskip

\subsection{2.1.2 Các ví dụ thực tế về Usability}
\begin{itemize}[leftmargin=1.5em]
\item \textbf{Dễ sử dụng (Ease of Use):}  
\end{itemize}
Một ứng dụng mua sắm trực tuyến cần có giao diện đơn giản, rõ ràng để người dùng tìm kiếm và mua hàng trong vài bước. Nút \textit{"Thêm vào giỏ hàng"} và \textit{"Thanh toán"} phải được đặt ở vị trí trực quan, nổi bật, dễ bấm trúng trên màn hình di động.

\begin{itemize}[leftmargin=1.5em]
\item \textbf{Hiệu suất tìm kiếm (Search Efficiency):}  
\end{itemize}
Ứng dụng đọc tin tức cần trang bị bộ công cụ tìm kiếm mạnh mẽ, cho phép người dùng truy cập bài viết tức thì theo từ khóa, danh mục chủ đề, tác giả hoặc mốc thời gian.

\begin{itemize}[leftmargin=1.5em]
\item \textbf{Dễ đọc và dễ hiểu (Readability \& Comprehension):}  
\end{itemize}
Trong các ứng dụng học trực tuyến, văn bản phải được định dạng theo cấu trúc phân cấp, cỡ chữ chuẩn, khoảng cách dòng hợp lý, tránh các khối chữ dày đặc gây mỏi mắt người học.


\begin{figure}[htbp]
\centering
\includegraphics[width=0.75\textwidth]{images/image34.png}
\caption{Các yếu tố cấu thành trải nghiệm người dùng}
\end{figure}
\begin{figure}[htbp]
\centering
\includegraphics[width=0.75\textwidth]{images/image35.png}
\caption{Ví dụ thực tế về tối ưu hóa tính khả dụng}
\end{figure}

\bigskip\hrule\bigskip

\subsection{2.1.3 Năm thành phần chất lượng cốt lõi của Usability (Theo Jakob Nielsen)}

Mô hình chất lượng Usability kinh điển của Jakob Nielsen bao gồm 5 khía cạnh then chốt:


\begin{verbatim}
                          ┌───────────────────────────┐
                          │   5 THÀNH PHẦN USABILITY  │
                          └─────────────┬─────────────┘
                                        │
        ┌───────────────┬───────────────┼───────────────┬───────────────┐
        ▼               ▼               ▼               ▼               ▼
  Learnability     Efficiency      Memorability       Errors       Satisfaction
   (Dễ học)        (Hiệu suất)      (Dễ nhớ)        (Xử lý lỗi)    (Sự hài lòng)
\end{verbatim}

\begin{enumerate}[leftmargin=1.5em]
\item \textbf{Learnability (Khả năng học hỏi / Sự dễ học):}
\end{enumerate}
\begin{itemize}[leftmargin=1.5em]
\item Đánh giá mức độ dễ dàng mà người dùng lần đầu tiên tiếp xúc với giao diện có thể thực hiện và hoàn thành các tác vụ cơ bản.
\item \textit{Ví dụ:} Một ứng dụng giao thông công cộng trên smartphone có giao diện tối giản, sử dụng các biểu tượng quốc tế quen thuộc (xe buýt, tàu điện, trạm dừng). Người dùng mới mở app lần đầu là biết ngay cách nhập điểm đi - điểm đến mà không cần bất kỳ tài liệu hướng dẫn nào.
\end{itemize}

\begin{figure}[htbp]
\centering
\includegraphics[width=0.75\textwidth]{images/image36.png}
\caption{Learnability - Sự dễ học trong thiết kế}
\end{figure}

\begin{enumerate}[leftmargin=1.5em]
\item \textbf{Efficiency (Hiệu suất sử dụng):}
\end{enumerate}
\begin{itemize}[leftmargin=1.5em]
\item Đo lường tốc độ và sự tiện lợi mà người dùng đã quen thuộc với hệ thống có thể hoàn thành công việc của họ.
\item \textit{Ví dụ:} Trong ứng dụng lập kế hoạch chuyến đi, thay vì phải gõ địa chỉ thủ công, người dùng có thể kéo thả trực tiếp các điểm tham quan từ bản đồ vào lịch trình theo từng khung giờ chỉ bằng một thao tác chuột.
\end{itemize}

\begin{figure}[htbp]
\centering
\includegraphics[width=0.75\textwidth]{images/image37.png}
\caption{Efficiency - Tối ưu hóa hiệu suất thao tác}
\end{figure}

\begin{enumerate}[leftmargin=1.5em]
\item \textbf{Memorability (Khả năng ghi nhớ):}
\end{enumerate}
\begin{itemize}[leftmargin=1.5em]
\item Đánh giá khả năng người dùng quay trở lại sử dụng sản phẩm sau một thời gian dài gián đoạn có thể phục hồi ngay sự thành thạo mà không phải học lại từ đầu hay không.
\item \textit{Bí quyết thiết kế:} Sử dụng tính quen thuộc (Familiarity) và tính nhất quán (Consistency) cao. Ví dụ, ứng dụng ghi chú duy trì các biểu tượng lưu, tạo mới, thùng rác chuẩn mực giúp người dùng dù nửa năm không mở app vẫn thao tác chính xác tức thì.
\end{itemize}
\begin{enumerate}[leftmargin=1.5em]
\item \textbf{Errors (Khả năng xử lý và phòng ngừa lỗi):}
\end{enumerate}
\begin{itemize}[leftmargin=1.5em]
\item Đo lường: Người dùng gặp phải bao nhiêu lỗi trong quá trình sử dụng? Các lỗi đó có nghiêm trọng không? Và hệ thống hỗ trợ họ nhận diện, khắc phục lỗi dễ dàng như thế nào?
\item \textit{Ví dụ:} Trong ứng dụng đặt hàng, nếu khách hàng nhập sai định dạng thẻ tín dụng, hệ thống phải hiển thị cảnh báo đỏ trực tiếp ngay dưới ô nhập liệu (Inline validation) và hướng dẫn rõ ràng: \textit{"Số thẻ gồm 16 chữ số, vui lòng kiểm tra lại"}, thay vì để khách hàng bấm gửi form rồi mới báo lỗi chung chung.
\end{itemize}

\begin{figure}[htbp]
\centering
\includegraphics[width=0.75\textwidth]{images/image38.png}
\caption{Errors - Ngăn ngừa và xử lý lỗi kịp thời}
\end{figure}
\begin{figure}[htbp]
\centering
\includegraphics[width=0.75\textwidth]{images/image39.png}
\caption{Hệ thống cảnh báo lỗi trực quan}
\end{figure}

\begin{enumerate}[leftmargin=1.5em]
\item \textbf{Satisfaction (Sự hài lòng):}
\end{enumerate}
\begin{itemize}[leftmargin=1.5em]
\item Đánh giá cảm xúc, sự thoải mái và mức độ hài lòng về mặt thẩm mỹ cũng như tiện ích mà người dùng cảm nhận được sau khi trải nghiệm sản phẩm.
\end{itemize}

\begin{figure}[htbp]
\centering
\includegraphics[width=0.75\textwidth]{images/image40.png}
\caption{Satisfaction - Mang lại cảm giác thỏa mãn cho người dùng}
\end{figure}

\bigskip\hrule\bigskip

\section{2.2 Các nguyên lý duyệt web của người dùng (User Web Browsing Principles)}

Nghiên cứu hành vi thực tế của người dùng trên môi trường Internet (theo Steve Krug và Nielsen Norman Group) chỉ ra 3 đặc điểm tâm lý cốt lõi:

\begin{enumerate}[leftmargin=1.5em]
\item \textbf{Người dùng không đọc kỹ từng từ – Họ chỉ "quét" (Scanning):}  
\end{enumerate}
Người dùng lướt mắt rất nhanh qua trang web để tìm kiếm các từ khóa nổi bật, tiêu đề in đậm, hình ảnh hoặc nút bấm liên quan trực tiếp đến nhu cầu hiện tại của họ. Do đó, văn bản trên web cần được chia thành các đoạn ngắn, có gạch đầu dòng và tiêu đề phụ rõ ràng.

\begin{enumerate}[leftmargin=1.5em]
\item \textbf{Quy tắc "Thỏa mãn vừa đủ" (Satisficing):}  
\end{enumerate}
Khác với lý thuyết kinh tế học cho rằng con người luôn cân nhắc kỹ lưỡng để chọn giải pháp tối ưu nhất, trên môi trường số, người dùng thường bấm vào lựa chọn đầu tiên có vẻ hợp lý và chấp nhận được. Họ không muốn mất thời gian so sánh toàn bộ danh sách.

\begin{enumerate}[leftmargin=1.5em]
\item \textbf{Người dùng "mò mẫm" qua giao diện (Muddle Through):}  
\end{enumerate}
Hầu như không ai đọc sách hướng dẫn sử dụng hay phần trợ giúp khi lướt web. Người dùng thích vừa mò mẫm vừa làm thử. Nếu bấm nhầm, họ chỉ mong muốn có nút Quay lại (Back) hoặc Hoàn tác (Undo) để làm lại từ đầu.


\bigskip\hrule\bigskip

\section{2.3 Phương pháp thiết kế lấy người dùng làm trung tâm (User-Centered Design - UCD)}

\subsection{2.3.1 Khái niệm UCD}
\begin{itemize}[leftmargin=1.5em]
\item \textbf{User-Centered Design (UCD):} Là phương pháp và triết lý thiết kế mà trong đó nhu cầu, mong muốn, bối cảnh và giới hạn của người dùng cuối được đặt làm trọng tâm ở mọi giai đoạn của quá trình phát triển sản phẩm.
\item Theo Tổ chức Thiết kế Tương tác (\textit{Interaction Design Foundation}): UCD là quy trình lặp đi lặp lại nhằm tối ưu hóa sản phẩm xoay quanh cách thức con người muốn và cần sử dụng nó.
\item \textbf{Lợi ích kinh tế:} UCD giúp giảm thiểu tối đa số lần phải đập đi xây lại hoặc sửa đổi mã nguồn tốn kém sau khi sản phẩm đã phát hành, đồng thời gia tăng đột biến sự gắn kết của khách hàng.
\end{itemize}

\begin{figure}[htbp]
\centering
\includegraphics[width=0.75\textwidth]{images/image41.png}
\caption{Quy trình tạo giao diện người dùng}
\end{figure}
\begin{figure}[htbp]
\centering
\includegraphics[width=0.75\textwidth]{images/image42.png}
\caption{Mô hình thiết kế lấy người dùng làm trung tâm (UCD)}
\end{figure}

\bigskip\hrule\bigskip

\subsection{2.3.2 Năm nguyên tắc thiết kế UCD cốt lõi}
\begin{enumerate}[leftmargin=1.5em]
\item \textbf{Cho người dùng tham gia ngay từ đầu:} Đưa người dùng thực tế vào quá trình đóng góp ý kiến ngay từ khâu lên ý tưởng sơ khởi.
\item \textbf{Đồng cảm sâu sắc với người dùng (Empathy):} Đặt mình vào vị trí, tâm lý và hoàn cảnh của người dùng để hiểu rõ những khó khăn thực tế mà họ đối mặt.
\item \textbf{Nhiều vòng lặp phản hồi (Iterative Feedback Loops):} Liên tục kiểm tra, thu thập ý kiến, sửa đổi và kiểm tra lại qua nhiều chu kỳ phát triển.
\item \textbf{Đánh giá dựa trên trải nghiệm thực tế:} Đánh giá sản phẩm trong bối cảnh thực tế mà người dùng sẽ sử dụng, không đánh giá dựa trên môi trường giả lập phi thực tế của phòng làm việc.
\item \textbf{Hợp tác liên chức năng:} Kết hợp góc nhìn của chuyên gia thiết kế, kỹ sư phần mềm, chuyên viên nghiệp vụ và dữ liệu phân tích.
\end{enumerate}

\bigskip\hrule\bigskip

\subsection{2.3.3 Sáu phương pháp nghiên cứu người dùng trong UCD}
\begin{enumerate}[leftmargin=1.5em]
\item \textbf{Focus Groups (Nhóm tập trung):} Thảo luận nhóm định hướng để thu thập ý kiến và thái độ.
\item \textbf{Questionnaires \& Surveys (Bảng câu hỏi \& Khảo sát):} Thu thập dữ liệu thống kê định lượng trên diện rộng.
\item \textbf{Interviews (Phỏng vấn cá nhân):} Khai thác sâu các động cơ, cảm xúc và thói quen cá nhân.
\item \textbf{Usability Testing (Kiểm thử tính dùng được):} Quan sát trực tiếp người dùng thực hiện các kịch bản tác vụ cụ thể trên sản phẩm.
\item \textbf{Card Sorting (Sắp xếp thẻ):} Giúp tổ chức và phân cấp cấu trúc thông tin khoa học.
\item \textbf{Participatory Design (Thiết kế có sự tham gia):} Mời người dùng cùng tham gia phác thảo giải pháp thiết kế.
\end{enumerate}

\begin{figure}[htbp]
\centering
\includegraphics[width=0.75\textwidth]{images/image45.png}
\caption{Sáu phương pháp nghiên cứu UCD phổ biến}
\end{figure}

\bigskip\hrule\bigskip

\subsection{2.3.4 Các giai đoạn trong quy trình UCD}
Một quy trình UCD tiêu chuẩn (từ 5 đến 7 giai đoạn) diễn ra tuần tự và lặp lại như sau:

\begin{enumerate}[leftmargin=1.5em]
\item \textbf{Hiểu và xác định người dùng:} Nghiên cứu bối cảnh, đặc điểm nhân khẩu học và mục tiêu của người dùng.
\item \textbf{Thu thập yêu cầu người dùng:} Phỏng vấn, khảo sát để lập danh sách yêu cầu chức năng và phi chức năng.
\item \textbf{Phân tích nhiệm vụ (Task Analysis):} Mô hình hóa các bước thao tác nghiệp vụ cần giải quyết.
\item \textbf{Card Sorting \& Kiến trúc thông tin:} Sắp xếp cấu trúc cây menu, liên kết điều hướng và sitemap.
\item \textbf{Thiết kế giải pháp sơ khởi (Conceptual Design \& Wireframing):} Vẽ khung xương bố cục các màn hình chính.
\item \textbf{Xây dựng nguyên mẫu thử nghiệm (Prototyping):} Tạo các bản mẫu tương tác từ mức độ thấp đến cao.
\item \textbf{Đánh giá \& Kiểm thử khả dụng (Usability Evaluation):} Cho người dùng thao tác thực tế, ghi nhận lỗi và tinh chỉnh trước khi phát hành.
\end{enumerate}

\bigskip\hrule\bigskip

\section{2.4 Tạo mẫu thử (Prototyping) và các mức độ trung thực của bản vẽ}

\subsection{2.4.1 Prototyping là gì? Vì sao bắt buộc phải làm Prototype?}
\begin{itemize}[leftmargin=1.5em]
\item \textbf{Prototyping:} Là quá trình thực nghiệm hóa các ý tưởng thiết kế trừu tượng thành các phiên bản mô phỏng hữu hình mà con người có thể quan sát, chạm vào và tương tác trực tiếp.
\item \textbf{Công cụ hỗ trợ chuyên dụng:} Figma, ProtoPie, Adobe XD, Axure RP...
\item \textbf{Tầm quan trọng sống còn:}
\item Prototype cho phép thử nghiệm sớm trên nhóm khách hàng mục tiêu ở quy mô nhỏ.
\item Chi phí sửa đổi một bản vẽ Prototype rẻ hơn gấp \textbf{100 lần} so với chi phí phải sửa đổi lại toàn bộ hệ thống phần mềm sau khi lập trình viên đã viết mã nguồn xong.
\item Prototype là công cụ giao tiếp hữu hiệu nhất giữa khách hàng, nhà thiết kế và đội ngũ lập trình, xóa bỏ mọi hiểu lầm về mặt nghiệp vụ.
\end{itemize}

\bigskip\hrule\bigskip

\subsection{2.4.2 So sánh Interactive Prototype vs Static Prototype}
\begin{itemize}[leftmargin=1.5em]
\item \textbf{Interactive Prototype (Clickable Prototype - Nguyên mẫu tương tác):}
\item Được xây dựng bằng các phần mềm chuyên nghiệp, cho phép người dùng click chuột, vuốt chạm trên điện thoại, kích hoạt các hiệu ứng chuyển trang, mở popup như một ứng dụng thật.
\item Phản ánh trung thực 100\% cảm giác sử dụng thực tế.
\item \textbf{Static Prototype (Nguyên mẫu tĩnh):}
\item Là tập hợp các bản vẽ màn hình đứng yên (trên giấy hoặc file ảnh).
\item Khi thử nghiệm, người điều phối sẽ tự tay thay đổi các tờ giấy hoặc đổi ảnh màn hình tiếp theo dựa trên thao tác chỉ tay của người dùng (kỹ thuật \textit{Wizard of Oz}).
\end{itemize}

\bigskip\hrule\bigskip

\subsection{2.4.3 Sự tiến hóa từ Wireframe, Mockup đến Prototype}
\begin{enumerate}[leftmargin=1.5em]
\item \textbf{Wireframe (Bản vẽ khung xương):}  
\end{enumerate}
\begin{itemize}[leftmargin=1.5em]
\item Bản vẽ thô sơ, thường chỉ dùng hai màu đen trắng hoặc xám.  
\item Tập trung định hình cấu trúc thông tin, bố cục khu vực, luồng người dùng (\textbf{User Flows}) mà không chứa hình ảnh hay màu sắc đồ họa chi tiết.
\end{itemize}
\begin{enumerate}[leftmargin=1.5em]
\item \textbf{Mockup (Bản vẽ trực quan tĩnh):}  
\end{enumerate}
\begin{itemize}[leftmargin=1.5em]
\item Đã được khoác lên màu sắc, kiểu chữ (typography), icon, logo và hình ảnh hoàn chỉnh.  
\item Như một căn nhà đã được sơn phết, trang trí nội thất 3D đẹp đẽ nhưng các cửa chưa thể mở được (chưa có tương tác động).
\end{itemize}
\begin{enumerate}[leftmargin=1.5em]
\item \textbf{Prototype (Nguyên mẫu tương tác hoàn chỉnh):}  
\end{enumerate}
\begin{itemize}[leftmargin=1.5em]
\item Kết hợp cấu trúc của Wireframe và thẩm mỹ của Mockup, bổ sung thêm linh hồn chuyển động và khả năng phản hồi tương tác sống động, chính xác tuyệt đối.
\end{itemize}

\bigskip\hrule\bigskip

\section{2.5 Các nguyên tắc thiết kế tương tác nâng cao và đặc điểm người dùng}

\subsection{2.5.1 Ba "Luật Vàng" trong thiết kế giao diện}
\textit{(Đúc kết bởi Ben Shneiderman và Theo Mandel)}


\begin{myquote}
1. \textbf{Người dùng dễ điều khiển (Place Users in Control):}
- Người dùng phải luôn cảm thấy mình là người làm chủ hệ thống, chứ không phải bị hệ thống ép buộc hay kiểm soát.
- Cho phép người dùng dễ dàng hủy bỏ, thoát ra khỏi tác vụ hoặc quay lại bước trước bất kỳ lúc nào.

2. \textbf{Người dùng ít phải nhớ (Reduce Users' Memory Load):}
- Giới hạn khả năng nhớ tức thời (trí nhớ ngắn hạn) của con người rất hạn hẹp: bộ não chỉ ghi nhớ được khoảng \textbf{$7 \pm 2$ mẩu thông tin} cùng một lúc.
- Hệ thống phải hiển thị thông tin rõ ràng trước mắt thay vì bắt người dùng phải nhớ từ màn hình này sang màn hình khác.
- Cung cấp các giá trị mặc định (defaults) thông minh và phím tắt (shortcuts) hỗ trợ thao tác nhanh.

3. \textbf{Giao diện toàn vẹn và nhất quán (Make the Interface Consistent):}
- Duy trì sự đồng bộ về ngôn ngữ, ký hiệu, vị trí các nút chức năng và luồng tương tác trên toàn bộ hệ thống phần mềm.
\end{myquote}

\bigskip\hrule\bigskip

\subsection{2.5.2 Tám tiêu chuẩn thiết kế tương tác UCD mở rộng}

\begin{enumerate}[leftmargin=1.5em]
\item \textbf{Tương thích (Compatibility):}
\end{enumerate}
\begin{itemize}[leftmargin=1.5em]
\item \textit{Tương thích người dùng:} Phù hợp với năng lực nhận thức và thói quen của người dùng.
\item \textit{Tương thích nhiệm vụ:} Hỗ trợ tối ưu để hoàn thành đúng mục tiêu công việc.
\end{itemize}
\begin{enumerate}[leftmargin=1.5em]
\item \textbf{Thẩm mỹ và Rõ ràng (Aesthetics \& Clarity):}
\end{enumerate}
\begin{itemize}[leftmargin=1.5em]
\item Giao diện có tính thẩm mỹ cuốn hút, thị giác cân đối, phân cấp thông tin rõ ràng, không gây rối mắt.
\end{itemize}
\begin{enumerate}[leftmargin=1.5em]
\item \textbf{Dễ cấu hình và Nhất quán (Configurability \& Consistency):}
\end{enumerate}
\begin{itemize}[leftmargin=1.5em]
\item Cho phép người dùng dễ dàng cá nhân hóa giao diện (đổi giao diện sáng/tối, sắp xếp bảng làm việc) nhưng vẫn tuân thủ tính nhất quán toàn hệ thống.
\end{itemize}
\begin{enumerate}[leftmargin=1.5em]
\item \textbf{Trực tiếp và Hiệu quả (Directness \& Efficiency):}
\end{enumerate}
\begin{itemize}[leftmargin=1.5em]
\item Cung cấp phương thức hoàn thành tác vụ trực tiếp nhất, giảm thiểu tối đa các bước trung gian không cần thiết.
\end{itemize}
\begin{enumerate}[leftmargin=1.5em]
\item \textbf{Thân thiện và Mềm dẻo (Friendliness \& Flexibility):}
\end{enumerate}
\begin{itemize}[leftmargin=1.5em]
\item Sử dụng ngôn từ ấm áp, tôn trọng người dùng. Chấp nhận nhiều cách thức thao tác khác nhau (dùng chuột, phím tắt hoặc cảm ứng).
\end{itemize}
\begin{enumerate}[leftmargin=1.5em]
\item \textbf{Bỏ qua lỗi và Khôi phục trạng thái (Forgiveness \& Recovery):}
\end{enumerate}
\begin{itemize}[leftmargin=1.5em]
\item Hệ thống có lòng bao dung (Forgiveness), dung thứ và bỏ qua một số lỗi bất cẩn thông thường của người dùng mà không làm gián đoạn công việc; luôn hỗ trợ khôi phục trạng thái an toàn.
\end{itemize}
\begin{enumerate}[leftmargin=1.5em]
\item \textbf{Trách nhiệm và Đơn giản (Responsiveness \& Simplicity):}
\end{enumerate}
\begin{itemize}[leftmargin=1.5em]
\item Hệ thống phải có trách nhiệm phản hồi tức thì mọi yêu cầu của người dùng; giữ cho cấu trúc giao diện luôn đơn giản, tinh giản tối đa.
\end{itemize}
\begin{enumerate}[leftmargin=1.5em]
\item \textbf{Trong suốt, Dễ hiểu, Dễ học (Transparency, Understandability, Learnability):}
\end{enumerate}
\begin{itemize}[leftmargin=1.5em]
\item Tính \textit{Trong suốt (Transparency)} cho phép người dùng tập trung hoàn toàn vào nội dung công việc của họ mà không bị phân tâm bởi bản thân giao diện phần mềm.
\end{itemize}

\begin{figure}[htbp]
\centering
\includegraphics[width=0.75\textwidth]{images/image43.png}
\caption{Các tiêu chuẩn thiết kế tương tác UCD}
\end{figure}
\begin{figure}[htbp]
\centering
\includegraphics[width=0.75\textwidth]{images/image44.png}
\caption{Hệ thống nguyên tắc trải nghiệm tương tác}
\end{figure}

\bigskip\hrule\bigskip

\section{2.6 Chi tiết các phương pháp nghiên cứu và kiểm thử UCD}

\subsection{2.6.1 Focus Groups (Nhóm tập trung)}
\begin{itemize}[leftmargin=1.5em]
\item \textbf{Khái niệm:} Một nhóm từ 6 đến 10 người dùng đại diện được mời đến tham gia một buổi thảo luận chuyên sâu kéo dài 1 - 2 tiếng, dưới sự dẫn dắt của một người điều phối (Moderator) chuyên nghiệp.
\item \textbf{Mục tiêu:} Thu thập những cảm nhận, suy nghĩ tự nhiên, thái độ và mong muốn tiềm ẩn đối với sản phẩm mới hoặc sản phẩm hiện tại.
\item \textbf{Yếu tố thành công:} Người điều phối phải giữ thái độ trung lập, khuyến khích mọi người cởi mở chia sẻ, không để một cá nhân nào áp đặt quan điểm lên toàn nhóm.
\end{itemize}

\bigskip\hrule\bigskip

\subsection{2.6.2 Usability Testing (Kiểm thử tính dùng được)}
\begin{itemize}[leftmargin=1.5em]
\item \textbf{Khái niệm:} Là phương pháp quan sát trực tiếp người dùng thao tác trên giao diện sản phẩm mẫu để đánh giá khả năng sử dụng và phát hiện các rào cản tương tác.
\item \textbf{Các câu hỏi cốt lõi khi tổ chức Usability Testing:}
\item \textit{Ai phải tham gia?} Người dùng thực tế đại diện cho Persona, cùng sự quan sát của Tester, UX Designer, Product Manager và Developers.
\item \textit{Cần bao nhiêu người thử nghiệm?} Theo nghiên cứu của Jakob Nielsen, chỉ cần kiểm thử với \textbf{5 người dùng} là đã có thể phát hiện được tới \textbf{85\% các lỗi khả dụng nghiêm trọng nhất} của sản phẩm. Một nhóm thử nghiệm từ 5 đến 10 người là tỷ lệ vàng giữa chi phí và hiệu quả.
\item \textit{Có cần ghi hình không?} Rất cần thiết. Việc ghi âm, quay màn hình và biểu cảm khuôn mặt giúp đội ngũ có bằng chứng khách quan, không thể chối cãi về những điểm khiến người dùng bối rối.
\end{itemize}

\subsubsection{Bài học kinh nghiệm từ Microsoft (Năm 1995):}
\begin{itemize}[leftmargin=1.5em]
\item Khi phát triển các tính năng mới cho bộ ứng dụng văn phòng, Microsoft đã đầu tư xây dựng các phòng thí nghiệm Usability hiện đại.
\item Mỗi tính năng mới đều được đưa vào phòng Lab kiểm nghiệm với nhóm 10 người dùng ngẫu nhiên, ghi hình toàn bộ thao tác và đo đạc thời gian hoàn thành.
\item Điển hình là bộ phần mềm \textbf{Microsoft Office 4.0} đã trải qua hơn \textbf{8.000 giờ thử nghiệm tính tiện dụng} trong phòng lab trước khi phát hành chính thức, tạo nên chuẩn mực về sự tiện ích của phần mềm văn phòng trên toàn thế giới.
\end{itemize}

\begin{figure}[htbp]
\centering
\includegraphics[width=0.75\textwidth]{images/image46.png}
\caption{Mô hình phòng thí nghiệm Usability Testing của các tập đoàn công nghệ}
\end{figure}
\begin{figure}[htbp]
\centering
\includegraphics[width=0.75\textwidth]{images/image47.png}
\caption{Bảng phân công và khối lượng công việc kiểm thử khả dụng}
\end{figure}

\bigskip\hrule\bigskip

\subsection{2.6.3 Card Sorting (Sắp xếp thẻ đánh giá)}
\begin{itemize}[leftmargin=1.5em]
\item \textbf{Khái niệm:} Cho phép người dùng sắp xếp các mẩu giấy ghi chú (Post-it notes) hoặc thẻ số đại diện cho các nội dung/tính năng vào từng nhóm có ý nghĩa với họ.
\item \textbf{Phân loại:}
\item \textit{Open Card Sorting (Sắp xếp mở):} Người dùng tự do nhóm các thẻ và tự đặt tên cho từng nhóm $
\end{itemize}
ightarrow$ Giúp khám phá tư duy phân loại tự nhiên của người dùng.

\begin{itemize}[leftmargin=1.5em]
\item \textit{Closed Card Sorting (Sắp xếp đóng):} Người dùng sắp xếp các thẻ vào các nhóm danh mục đã được định nghĩa sẵn $
\end{itemize}
ightarrow$ Giúp kiểm tra độ chính xác của cấu trúc hiện có.

\begin{itemize}[leftmargin=1.5em]
\item \textbf{Ứng dụng:} Thiết kế kiến trúc thông tin, cấu trúc cây danh mục hàng hóa trên sàn thương mại điện tử và hệ thống menu sitemap.
\end{itemize}

\bigskip\hrule\bigskip

\subsection{2.6.4 Participatory Design (Thiết kế có sự tham gia)}
\begin{itemize}[leftmargin=1.5em]
\item Mời người dùng đóng vai trò như một thành viên trong nhóm thiết kế, trực tiếp tham gia vẽ ý tưởng, thảo luận giải pháp và ra quyết định thiết kế.
\item Giúp xóa bỏ hoàn toàn khoảng cách giữa suy nghĩ chủ quan của kỹ sư và mong muốn thực tế của người dùng.
\end{itemize}

\bigskip\hrule\bigskip

\subsection{2.6.5 Questionnaires \& Surveys (Bảng câu hỏi \& Khảo sát)}
\begin{itemize}[leftmargin=1.5em]
\item Thiết lập hệ thống câu hỏi trắc nghiệm hoặc thang đo Likert (1 đến 5 điểm) để thu thập dữ liệu từ hàng trăm, hàng ngàn người dùng từ xa qua Internet.
\item Cung cấp các số liệu thống kê đáng tin cậy phục vụ cho việc ra quyết định ở quy mô lớn.
\end{itemize}

\bigskip\hrule\bigskip

\subsection{2.6.6 Interviews (Phỏng vấn cá nhân sâu)}
\begin{itemize}[leftmargin=1.5em]
\item Tiến hành phỏng vấn 1-1 giữa chuyên viên nghiên cứu và người dùng.
\item \textbf{Ưu điểm vượt trội:} Đi sâu vào từng chi tiết ngóc ngách tâm lý, giải thích ngay các hiểu lầm giữa hai bên và thu thập được các dữ liệu định tính vô giá.
\end{itemize}

\bigskip\hrule\bigskip

\section{2.7 Case Study thực hành: Thiết kế hệ thống đặt pizza trực tuyến (Online Pizza Ordering)}

Một dự án thiết kế hoàn chỉnh theo quy trình UCD kinh điển:


\subsection{2.7.1 Khảo sát thực tế và Xác định yêu cầu}
\begin{itemize}[leftmargin=1.5em]
\item \textbf{Mục tiêu:} Thiết kế hệ thống đặt bánh pizza trực tuyến thay thế hoặc bổ trợ cho hai kênh truyền thống: Mua tại quầy và Gọi điện thoại.
\item \textbf{Đề xuất các tính năng mới:} Phiếu giảm giá trực tuyến, đặt hàng trước theo giờ hẹn, khuyến nghị món ăn thông minh, hệ thống đánh giá của cộng đồng, so sánh giá.
\item \textbf{Khảo sát thực địa:}
\item Đến tận cửa hàng pizza để quan sát quy trình làm việc thực tế: quan sát khách mua hàng trực tiếp, quan sát nhân viên thu ngân nhập máy POS.
\item Phỏng vấn sâu nhân viên và ít nhất 3 khách hàng với các câu hỏi: Họ thích tự phối nhân bánh/đế bánh hay chọn bánh có sẵn trong menu? Khách hàng có mong muốn hỗ trợ trực tuyến không?
\end{itemize}

\bigskip\hrule\bigskip

\subsection{2.7.2 Phân tích nhiệm vụ (Task Analysis): So sánh quy trình hiện tại}

Dựa trên kết quả phỏng vấn, lập bảng ma trận phân tích 5 nhiệm vụ cốt lõi:


\begin{table}[htbp]
\centering
\begin{tabularx}{\textwidth}{p{0.2\textwidth} p{0.38\textwidth} p{0.38\textwidth}}
\toprule
\textbf{Bước nhiệm vụ} & \textbf{Kênh Tại Quầy (In-Store)} & \textbf{Kênh Qua Điện Thoại (Phone Ordering)} \\
\midrule
\textbf{1. Khởi tạo giao dịch} & Khách trực tiếp đến cửa hàng, đứng trước quầy thu ngân. & Khách quay số gọi điện đến đường dây nóng của cửa hàng. \\
\textbf{2. Xem món ăn} & Xem trực tiếp thực đơn giấy, ngắm ảnh minh họa trên bảng LED. & Xem thực đơn đã có từ trước, hỏi nhân viên tư vấn món ăn ngon. \\
\textbf{3. Chọn món} & Lựa chọn loại đế bánh, cỡ bánh, nhân bánh, món ăn phụ và nước uống. & Nhân viên ghi chép lựa chọn đế bánh, nhân bánh, món phụ và nước uống. \\
\bottomrule
\end{tabularx}
\end{table}
| \textbf{4. Xác nhận \& Thanh toán} | Xác nhận lại đơn hàng $

ightarrow$ Trả tiền mặt hoặc quẹt thẻ tại quầy. (Nếu không ưng ý, quay lại bước 2-3). | Xác nhận địa chỉ nhận hàng $

ightarrow$ Cam kết trả tiền mặt khi nhận hàng (COD). |


\bigskip\hrule\bigskip

\subsection{2.7.3 Sàng lọc 5 chức năng then chốt cho hệ thống Online}
Trong số các ý tưởng đề xuất, dựa vào dữ liệu phỏng vấn người dùng thực tế, nhóm thiết kế quyết định \textbf{chọn 4 chức năng cốt lõi và loại bỏ 1 chức năng} như sau:

\begin{enumerate}[leftmargin=1.5em]
\item \textbf{Xem trước ảnh món ăn chất lượng cao:} Khách hàng muốn nhìn thấy rõ chiếc bánh trước khi đặt. \textit{(Được chọn)}
\item \textbf{Xem trước giá cả niêm yết minh bạch:} Tính toán tổng chi phí tức thì bao gồm phí ship. \textit{(Được chọn)}
\item \textbf{Xem thông tin về lượng Calorie tương ứng:} Khách hàng hiện đại rất quan tâm đến sức khỏe và chế độ dinh dưỡng. \textit{(Được chọn)}
\item \textbf{Tùy biến thành phần nguyên liệu (Customization):} Khảo sát cho thấy phần lớn khách hàng muốn tự tay chọn đế dày/mỏng, viền phô mai và thêm bớt nhân theo khẩu vị riêng. \textit{(Được chọn)}
\item \textbf{So sánh giá cả với các thương hiệu pizza khác:} Khảo sát cho thấy hương vị bánh và tốc độ giao hàng là yếu tố quyết định hàng đầu, khách hàng không muốn xem quảng cáo so sánh giá của đối thủ trên chính web của quán. \textit{(LOẠI BỎ)}
\end{enumerate}

\bigskip\hrule\bigskip

\subsection{2.7.4 Xây dựng 3 Kịch bản người dùng điển hình (Scenarios)}
\begin{itemize}[leftmargin=1.5em]
\item \textbf{Kịch bản 1 (Khách hàng thân thiết):} Người thường xuyên ăn pizza, đặt hàng nhanh tại nhà vào giờ ăn tối, cần giao hàng gấp trong vòng 30 phút.
\item \textbf{Kịch bản 2 (Khách hàng mới):} Lần đầu tiên đặt bánh qua web, phân vân chưa biết chọn loại nào, cần các gợi ý "Best-seller" và khuyến mãi chào mừng.
\item \textbf{Kịch bản 3 (Khách hàng đặt trước cho sự kiện):} Khách hàng đặt mua tại công sở, đặt trước 2 tiếng cho buổi liên hoan công ty vào lúc 16h00.
\end{itemize}

\bigskip\hrule\bigskip

\subsection{2.7.5 Thiết kế các màn hình tương tác chính}
\begin{itemize}[leftmargin=1.5em]
\item \textbf{Màn hình chọn món chính (Pizza Selection):} Lựa chọn kích thước (S/M/L), loại đế (dày/mỏng/giòn), viền nhân phô mai/xúc xích và topping.
\item \textbf{Màn hình chọn món phụ \& Đồ uống (Sides \& Drinks):} Gợi ý mì Ý, gà rán, salad, coca-cola.
\item \textbf{Màn hình Xác nhận đơn hàng \& Thanh toán:} Tóm tắt đơn hàng, nhập mã voucher giảm giá, điền địa chỉ giao hàng, chọn thời gian giao và thanh toán online.
\item \textbf{Kiểm thử kịch bản (Walk-through):} Cho người dùng thao tác mẫu thử để kiểm tra xem luồng đặt hàng có trơn tru hay không trước khi chuyển sang giai đoạn lập trình.
\end{itemize}

\bigskip\hrule\bigskip

\subsection{2.7.6 Các chủ đề bài tập lớn thực hành (Gợi ý đồ án môn học)}
\begin{itemize}[leftmargin=1.5em]
\item \textbf{Chủ đề 1:} Hệ thống đặt bàn và gọi món trực tuyến cho chuỗi nhà hàng ẩm thực.
\item \textbf{Chủ đề 2:} Ứng dụng điểm danh sinh viên thông minh ứng dụng nhận diện khuôn mặt.
\item \textbf{Chủ đề 3:} Cổng thông tin và quản lý hoạt động câu lạc bộ sinh viên trường đại học.
\end{itemize}


\bigskip\hrule\bigskip

\chapter{CHƯƠNG 3: KỸ THUẬT THIẾT KẾ CÁC THÀNH PHẦN TRÊN GIAO DIỆN}

\section{3.1 Màn hình và nguyên tắc thiết kế màn hình}

\subsection{3.1.1 Bản chất của thiết kế màn hình}
\begin{itemize}[leftmargin=1.5em]
\item Thiết kế giao diện người dùng thực chất là nghệ thuật tổ chức và trình bày tất cả những gì hiển thị trên màn hình.
\item Giao diện là những gì trực tiếp đập vào mắt người sử dụng $\rightarrow$ Mục tiêu tối thượng là phải thiết kế sao cho \textbf{người dùng nhìn vào là hiểu ngay lập tức} ý nghĩa và chức năng của từng thành phần.
\item Luôn ưu tiên thiết kế cho màn hình cơ sở có độ phân giải chuẩn, sắc nét và khả năng co giãn linh hoạt trên đa dạng kích thước hiển thị.
\end{itemize}

\bigskip\hrule\bigskip

\subsection{3.1.2 Quy trình 3 bước cốt lõi: Hỏi – Nghĩ – Thiết kế (Ask – Think – Design)}
\begin{enumerate}[leftmargin=1.5em]
\item \textbf{Hỏi (Ask):}  
\end{enumerate}
\begin{itemize}[leftmargin=1.5em]
\item Người dùng đang làm gì?  
\item Họ truy cập vào màn hình này trong bối cảnh nào và với tâm trạng ra sao?
\end{itemize}
\begin{enumerate}[leftmargin=1.5em]
\item \textbf{Nghĩ (Think):}  
\end{enumerate}
\begin{itemize}[leftmargin=1.5em]
\item Thông tin gì bắt buộc phải được hiển thị trước mắt họ?  
\item Người dùng cần thực hiện những phép so sánh hoặc đối chiếu dữ liệu nào?  
\item Thứ tự logic của các thao tác cần được diễn ra như thế nào?
\end{itemize}
\begin{enumerate}[leftmargin=1.5em]
\item \textbf{Thiết kế (Design):}  
\end{enumerate}
\begin{itemize}[leftmargin=1.5em]
\item \textbf{Thiết kế theo đúng chức năng (Form follows function).}  
\item Các tương tác người dùng đòi hỏi sẽ đóng vai trò quyết định cấu trúc bố cục (\textbf{Layout}) của màn hình.
\end{itemize}

\bigskip\hrule\bigskip

\section{3.2 Bố cục và hệ thống lưới (Layout \& Grid System)}

\subsection{3.2.1 Vai trò của hệ thống lưới trong UI Design}
\begin{itemize}[leftmargin=1.5em]
\item \textbf{Lưới (Grid):} Là một cấu trúc khung vô hình bao gồm các đường gióng dọc và ngang, giúp phân chia không gian và giữ vững trật tự bố cục của toàn bộ thiết kế.
\item Lưới được ví như \textbf{"xương sống"} của mỗi bản thiết kế. Một bản thiết kế có cấu trúc tốt sẽ sở hữu hệ thống lưới phân cấp rõ ràng, tạo nên sự cân đối, hài hòa thị giác và dễ dàng sử dụng.
\item Thiết lập hệ thống lưới là công việc đầu tiên và quan trọng nhất khi bắt tay vào thiết kế giao diện, đảm bảo tính nhất quán trên hàng trăm màn hình của dự án.
\end{itemize}

\bigskip\hrule\bigskip

\subsection{3.2.2 Các kiểu lưới cơ bản}
\begin{itemize}[leftmargin=1.5em]
\item \textbf{Lưới ngang (Defining Horizontal Layout):} Bao gồm các cột (\textbf{Columns}), khoảng trống giữa các cột (\textbf{Gutters}) và lề hai bên (\textbf{Margins}). Chịu trách nhiệm căn chỉnh các thành phần giao diện theo chiều ngang của màn hình (hệ thống 12 cột kinh điển trên Web).
\item \textbf{Lưới dọc (Defining Vertical Layout):} Là biến thể giúp thiết lập độ cao của các phần tử, phân chia các khối nội dung và quản lý khoảng trắng theo trục thẳng đứng.
\end{itemize}

\bigskip\hrule\bigskip

\subsection{3.2.3 Lưới linh động (Soft Grid) và Phương pháp căn chỉnh "Hình vuông đỏ"}
\begin{itemize}[leftmargin=1.5em]
\item \textbf{Hạn chế của lưới cứng:} Một hệ thống lưới sắp xếp quá cứng nhắc sẽ làm cho giao diện trở nên đơn điệu, rập khuôn và gây nhàm chán cho người dùng.
\item \textbf{Lưới linh động (Soft Grid):} Là giải pháp căn chỉnh các thành phần trên giao diện một cách nhịp nhàng, thanh thoát và tự nhiên hơn dựa trên một đơn vị cơ sở chuẩn.
\item \textbf{Quy tắc cơ sở (Base Unit):} Mọi khoảng cách lề (margin), khoảng cách đệm (padding), kích thước biểu tượng và chiều cao thành phần đều phải là \textbf{bội số của một số cơ sở} (phổ biến nhất trong thiết kế hiện đại là hệ số \textbf{8pt} hoặc \textbf{4pt}: 4, 8, 12, 16, 24, 32, 40, 48, 64px...).
\end{itemize}

\subsubsection{Quy trình thực hiện phương pháp "Hình vuông đỏ":}
\begin{enumerate}[leftmargin=1.5em]
\item \textbf{Bước 1:} Chọn con số cơ sở chuẩn cho toàn bộ dự án (ví dụ chuẩn: 8px).
\item \textbf{Bước 2:} Tạo một đối tượng hình vuông màu đỏ có kích thước cạnh đúng bằng con số cơ sở đã chọn ($8 \times 8$ px).
\item \textbf{Bước 3:} Sử dụng hình vuông màu đỏ làm thước đo chuẩn để căn chỉnh khoảng cách giữa các khối đối tượng: khoảng cách 1 ô (8px), 2 ô (16px), 3 ô (24px)... nhằm đảm bảo nhịp điệu thị giác nhất quán 100\%.
\end{enumerate}

\bigskip\hrule\bigskip

\section{3.3 Thứ tự các nhóm và thành phần giao diện}

\subsection{3.3.1 Thứ tự tự nhiên đối với người dùng}
\begin{itemize}[leftmargin=1.5em]
\item Thiết kế phải tuân thủ thứ tự nhận thức tự nhiên của người dùng: từ trên xuống dưới, từ trái sang phải theo thói quen đọc văn bản truyền thống.
\item Các thông tin quan trọng nhất, các tiêu đề bao quát phải được đặt ở góc trên bên trái; các nút hành động xác nhận cuối cùng (Submit, Hoàn tất) đặt ở phía dưới cùng bên phải.
\end{itemize}

\subsection{3.3.2 Thứ tự phím Tab (Tab Index)}
\begin{itemize}[leftmargin=1.5em]
\item Đối với biểu mẫu nhập liệu (Form) hoặc hộp thoại (Dialogue Box), việc thiết lập thứ tự di chuyển khi nhấn phím \textbf{Tab} là cực kỳ quan trọng.
\item Phím Tab phải nhảy tuần tự từ trường nhập liệu này sang trường nhập liệu kế tiếp một cách logic. Tuyệt đối không để con trỏ nhảy lung tung, gây ức chế cho người dùng nhập liệu bằng bàn phím.
\end{itemize}

\subsection{3.3.3 Thứ tự quy chuẩn đặc biệt}
\begin{itemize}[leftmargin=1.5em]
\item Trong các danh sách tra cứu lớn (danh bạ điện thoại, danh mục quốc gia, bảng từ điển), bắt buộc phải sắp xếp theo thứ tự bảng chữ cái alphabet (A $\rightarrow$ Z).
\end{itemize}

\bigskip\hrule\bigskip

\section{3.4 Kỹ thuật căn chỉnh dữ liệu (Data Alignment)}

\subsection{3.4.1 Căn chỉnh văn bản (Text Alignment)}
\begin{itemize}[leftmargin=1.5em]
\item \textbf{Quy luật thị giác:} Mắt người đọc văn bản theo hàng ngang từ trái sang phải.
\item \textbf{Căn lề trái (Left-aligned):} Là phương thức sắp xếp tự nhiên, dễ đọc nhất đối với mắt người, tạo ra một đường gióng thẳng đứng bên trái giúp mắt dễ dàng quay lại đầu dòng kế tiếp.
\item \textbf{Căn lề phải (Right-aligned):} Gây mỏi mắt nếu áp dụng cho đoạn văn bản dài vì điểm bắt đầu của mỗi dòng liên tục thay đổi. Chỉ nên dùng căn phải cho các nhãn chú thích ngắn hoặc đối sánh dữ liệu.
\item \textbf{Căn giữa (Center-aligned):} Chỉ sử dụng cho các tiêu đề ngắn, câu danh ngôn hoặc thông báo nổi bật dưới 3 dòng.
\end{itemize}

\bigskip\hrule\bigskip

\subsection{3.4.2 Căn chỉnh Tên riêng (Names Alignment)}
\begin{itemize}[leftmargin=1.5em]
\item \textbf{Quy tắc bắt buộc:} Luôn tách rời hai trường thông tin: \textbf{Họ} (Last Name) và \textbf{Tên} (First Name).
\item \textbf{Mục đích:} Phục vụ cho các thao tác nghiệp vụ thiết yếu như lọc dữ liệu, sắp xếp danh sách theo thứ tự bảng chữ cái của Tên và tối ưu hóa công cụ tìm kiếm người dùng.
\end{itemize}

\bigskip\hrule\bigskip

\subsection{3.4.3 Căn chỉnh Số liệu (Numbers Alignment)}
Số liệu trong thiết kế giao diện đòi hỏi sự khắt khe cao về mặt thị giác:

\begin{itemize}[leftmargin=1.5em]
\item \textbf{Nguyên tắc trực quan:} Chuỗi số dài hơn biểu thị một giá trị lớn hơn.
\item \textbf{Quy chuẩn căn lề:}
\item \textbf{Số nguyên (Integers):} Luôn luôn \textbf{căn lề bên phải (Right-aligned)}. Điều này giúp hàng đơn vị, hàng chục, hàng trăm, hàng nghìn thẳng tắp theo một đường gióng dọc, cho phép người dùng lướt mắt so sánh độ lớn của các con số ngay lập tức.
\item \textbf{Số thực / Số thập phân (Decimals):} Bắt buộc phải \textbf{căn thẳng hàng theo dấu thập phân} (dấu chấm \texttt{.} hoặc dấu phẩy \texttt{,}).
\item \textbf{Dữ liệu văn bản đi kèm:} Tiêu đề cột số cũng nên căn phải theo số liệu bên dưới để tạo sự liền mạch.
\item \textbf{Quy tắc phân đoạn ghi nhớ (Chunking):}  
\end{itemize}
Bộ não con người rất khó ghi nhớ và đọc chính xác một dãy số dài dằng dặc. Do đó, bắt buộc phải phân chia chuỗi số bằng dấu phân cách:

\begin{itemize}[leftmargin=1.5em]
\item \textit{Số tiền tệ:} Sử dụng dấu phân cách hàng nghìn (ví dụ: \texttt{2,502,560} thay vì \texttt{2502560}).
\item \textit{Số điện thoại / Số thẻ:} Chia thành các cụm 3 đến 4 chữ số (ví dụ: \texttt{0903 123 456} hoặc \texttt{4111 2222 3333 4444}).
\item \textbf{Dạng bảng nhiều hàng nhiều cột (Data Tables):}
\item Sử dụng các đường dẫn dắt (Leader lines) hoặc áp dụng màu sọc xen kẽ nhẹ nhàng (\textbf{Zebra striping}) giữa các dòng để dẫn mắt người dùng duyệt ngang qua các cột mà không bị lệch dòng.
\end{itemize}

\begin{figure}[htbp]
\centering
\includegraphics[width=0.75\textwidth]{images/image48.png}
\caption{Quy chuẩn căn chỉnh số liệu thẳng hàng dấu thập phân}
\end{figure}
\begin{figure}[htbp]
\centering
\includegraphics[width=0.75\textwidth]{images/image49.png}
\caption{Thiết kế bảng số liệu chuẩn mực và trực quan}
\end{figure}

\bigskip\hrule\bigskip

\section{3.5 Khoảng trắng (White Space) và các định luật Gestalt}

\subsection{3.5.1 Sức mạnh của Khoảng trắng trong thiết kế}
\begin{itemize}[leftmargin=1.5em]
\item \textbf{Khoảng trắng (White Space / Negative Space):} Không phải là "vùng chết" hay không gian lãng phí, mà là một trong những công cụ thiết kế quyền lực nhất.
\item Khoảng trắng tạo không gian cho mắt người nghỉ ngơi, giảm bớt căng thẳng thị giác, phân tách các nhóm thông tin và làm nổi bật các thành phần trọng tâm của sản phẩm.
\end{itemize}

\bigskip\hrule\bigskip

\subsection{3.5.2 Các định luật cảm nhận thị giác Gestalt}
\begin{itemize}[leftmargin=1.5em]
\item \textbf{Luật Cận kề (Law of Proximity):} Các đối tượng nằm gần nhau trong không gian sẽ được não bộ tự động nhóm lại và nhận thức như một thể thống nhất có liên quan mật thiết với nhau.
\item \textbf{Luật Tương tự (Law of Similarity):} Những thành phần có chung màu sắc, hình dáng hoặc kích thước sẽ được người dùng ngầm hiểu là có cùng vai trò hoặc cấp độ chức năng.
\end{itemize}

\bigskip\hrule\bigskip

\subsection{3.5.3 Phân loại khoảng trắng}
\begin{enumerate}[leftmargin=1.5em]
\item \textbf{Khoảng trắng trong trình bày văn bản (Micro White Space):}
\end{enumerate}
\begin{itemize}[leftmargin=1.5em]
\item Khoảng cách giữa các ký tự (\textit{Kerning}).
\item Khoảng cách giữa các từ (\textit{Tracking}).
\item Khoảng cách giữa các dòng (\textit{Line-height / Leading}): Tiêu chuẩn vàng là độ cao dòng bằng \textbf{140\% – 160\%} kích thước cỡ chữ.
\item Khoảng cách giữa các đoạn văn (\textit{Paragraph spacing}).
\end{itemize}
\begin{enumerate}[leftmargin=1.5em]
\item \textbf{Khoảng trắng trong bố trí giao diện (Macro White Space):}
\end{enumerate}
\begin{itemize}[leftmargin=1.5em]
\item Khoảng cách đệm (\textit{Padding}) bên trong các thẻ card, nút bấm.
\item Khoảng cách lề (\textit{Margin}) giữa các khối chức năng khác nhau.
\item Giúp người dùng phân biệt ranh giới giữa các khối dữ liệu mà không cần phải lạm dụng quá nhiều đường kẻ phân cách thô cứng.
\end{itemize}

\begin{figure}[htbp]
\centering
\includegraphics[width=0.75\textwidth]{images/image50.png}
\caption{Khoảng trắng trong thiết kế giao diện}
\end{figure}
\begin{figure}[htbp]
\centering
\includegraphics[width=0.75\textwidth]{images/image51.png}
\caption{Ứng dụng các quy luật Gestalt trong bố trí khoảng trắng}
\end{figure}

\bigskip\hrule\bigskip

\section{3.6 Nghệ thuật sử dụng màu sắc (Color) trong thiết kế UI}

\subsection{3.6.1 Tầm quan trọng và Yếu tố văn hóa của màu sắc}
\begin{itemize}[leftmargin=1.5em]
\item Màu sắc là yếu tố đầu tiên kích thích thị giác và khơi gợi cảm xúc của người dùng trước khi họ kịp đọc câu chữ.
\item \textbf{Lưu ý về văn hóa (Cultural Context):} Ý nghĩa của màu sắc thay đổi đáng kể tùy theo từng nền văn hóa:
\item \textit{Màu đỏ:} Ở các nước phương Tây thường mang nghĩa nguy hiểm, cảnh báo, dừng lại; nhưng tại Trung Quốc và một số nước châu Á lại biểu trưng cho may mắn, tài lộc và hạnh phúc.
\item \textit{Màu trắng:} Phương Tây coi là sự tinh khiết, thanh lịch; nhưng ở một số quốc gia phương Đông lại gợi nhắc sự mất mát, tang tóc.
\end{itemize}

\begin{figure}[htbp]
\centering
\includegraphics[width=0.75\textwidth]{images/image52.png}
\caption{Ý nghĩa màu sắc và sự khác biệt văn hóa}
\end{figure}

\bigskip\hrule\bigskip

\subsection{3.6.2 Ý nghĩa và ứng dụng của ba gam màu chủ đạo trong UI}
\begin{enumerate}[leftmargin=1.5em]
\item \textbf{Màu xanh dương (Blue):}
\end{enumerate}
\begin{itemize}[leftmargin=1.5em]
\item Là gam màu được ưa chuộng và sử dụng phổ biến nhất trong các sản phẩm kỹ thuật số (Facebook, Twitter, Microsoft, IBM, PayPal, Telegram).
\item \textbf{Tâm lý học:} Gợi lên cảm giác tin cậy, an toàn, ổn định, thông tuệ và chuyên nghiệp.
\item \textbf{Ứng dụng:} Lý tưởng cho các ứng dụng ngân hàng, tài chính số, y tế, mạng xã hội và phần mềm doanh nghiệp.
\end{itemize}
\begin{enumerate}[leftmargin=1.5em]
\item \textbf{Màu xanh lá cây (Green):}
\end{enumerate}
\begin{itemize}[leftmargin=1.5em]
\item Là màu sắc mà mắt người có độ nhạy cảm thị giác cao nhất trong dải quang phổ.
\item \textbf{Tâm lý học:} Biểu thị sự sinh sôi, thiên nhiên, hòa bình, tươi mới và an toàn.
\item \textbf{Ứng dụng:} Sử dụng cho các trạng thái thành công, nút Xác nhận hoàn tất, thanh toán an toàn, hoặc các ứng dụng nông nghiệp sạch, môi trường.
\end{itemize}
\begin{enumerate}[leftmargin=1.5em]
\item \textbf{Màu đỏ (Red):}
\end{enumerate}
\begin{itemize}[leftmargin=1.5em]
\item Màu sắc kích thích xung thần kinh mạnh mẽ nhất, thu hút ánh nhìn tức thì và kích thích nhịp tim.
\item \textbf{Tâm lý học:} Tình yêu, đam mê, sự khẩn cấp và nguy hiểm.
\item \textbf{Ứng dụng:} Tuyệt đối không dùng tràn lan. Chỉ dành riêng cho các cảnh báo nghiêm trọng, thông báo lỗi hệ thống, nút Xóa dữ liệu vĩnh viễn, hoặc huy hiệu thông báo chưa đọc (\textbf{Notification Badge}).
\end{itemize}

\bigskip\hrule\bigskip

\subsection{3.6.3 Ví dụ thiết kế bảng điều khiển vật lý (Physical Controls)}
\begin{itemize}[leftmargin=1.5em]
\item Xem xét bảng điều khiển của một lò nấu đa năng: Các nút bấm cần được nhóm lại theo đúng trình tự vận hành logic:
\end{itemize}
\[ \text{1. Loại làm nóng} \longrightarrow \text{2. Nhiệt độ} \longrightarrow \text{3. Thời gian nấu} \longrightarrow \text{4. Bắt đầu (Start)} \]
\begin{itemize}[leftmargin=1.5em]
\item Sử dụng màu sắc khác nhau để phân định rõ ràng ranh giới chức năng; vẽ đường viền bao quanh các nhóm nút có liên quan nhằm giảm thiểu tối đa sự nhầm lẫn của người sử dụng.
\end{itemize}

\bigskip\hrule\bigskip

\section{3.7 Kiểu chữ (Typography) và độ tương phản thị giác}

\begin{itemize}[leftmargin=1.5em]
\item \textbf{Kiểu chữ (Typography):} Truyền tải linh hồn, phong cách và thông điệp của toàn bộ sản phẩm đến tâm trí người dùng.
\item \textbf{Kích thước tiêu chuẩn:}
\item Nội dung thông thường (\textit{Body Text}): Sử dụng kích thước \textbf{10 – 12 pt} trên tài liệu in ấn (hoặc \textbf{14 – 16 px} trên màn hình số).
\item Tiêu đề cấp 1 (\textit{Heading 1}): $28 - 36$ px; Tiêu đề cấp 2 (\textit{Heading 2}): $20 - 24$ px.
\item \textbf{Quy tắc định dạng:}
\item Hạn chế lạm dụng in đậm (\textbf{Bold}) hoặc in nghiêng (\textit{Italic}) trong toàn bộ khối văn bản dài. Chỉ in đậm các từ khóa thực sự quan trọng cần ghi nhớ.
\item Tuyệt đối tránh sử dụng các font chữ uốn lượn, kỳ quái hoặc đã lỗi thời. Ưu tiên các font Sans-serif (không chân) hiện đại, rõ nét trên màn hình số như \textit{Inter, Roboto, San Francisco, Segoe UI}.
\item \textbf{Độ tương phản (Contrast Ratio):} Phải đảm bảo tỷ lệ tương phản tối thiểu \textbf{4.5:1} giữa chữ và màu nền theo tiêu chuẩn trợ năng quốc tế WCAG.
\end{itemize}

\begin{figure}[htbp]
\centering
\includegraphics[width=0.75\textwidth]{images/image53.png}
\caption{Typography - Lựa chọn kiểu chữ và phân cấp thị giác}
\end{figure}

\bigskip\hrule\bigskip

\section{3.8 Liên kết văn bản (Text Links)}

\begin{itemize}[leftmargin=1.5em]
\item Liên kết văn bản là phương tiện điều hướng căn bản nhất của mạng Internet.
\item \textbf{Quy ước bất thành văn:} Chữ màu xanh dương và có gạch chân bên dưới (\textit{Underlined}) là dấu hiệu nhận diện quốc tế kinh điển báo cho người dùng biết văn bản đó có thể nhấp chuột vào.
\item \textbf{Bốn trạng thái bắt buộc của Link:}
\item \textit{Link thường (Normal / Unvisited):} Màu xanh dương.
\item \textit{Rê chuột (Hover):} Đổi sắc thái màu hoặc xuất hiện gạch chân để báo hiệu sẵn sàng click.
\item \textit{Đang click (Active):} Phản hồi trạng thái nhấn.
\item \textit{Đã ghé thăm (Visited):} Chuyển sang tông màu tím nhẹ để người dùng không bấm lại trang cũ một cách vô ích.
\end{itemize}

\bigskip\hrule\bigskip

\section{3.9 Thiết kế Biểu tượng (Icons)}

\begin{itemize}[leftmargin=1.5em]
\item \textbf{Biểu tượng (Icon):} Là một chữ tượng hình nhỏ đại diện cho một chức năng, khái niệm hoặc trạng thái của hệ thống.
\item \textbf{Bốn nguyên tắc thiết kế Icon chuẩn mực:}
\end{itemize}
\begin{enumerate}[leftmargin=1.5em]
\item \textbf{Tính nhất quán (Consistency):} Toàn bộ bộ icon phải thống nhất về kích thước khung (Grid $24 \times 24$ px), độ dày đường nét (Stroke 2px), góc bo và phong cách (hoàn toàn là nét Outline hoặc hoàn toàn là mảng Solid).
\item \textbf{Tính quen thuộc (Familiarity):} Sử dụng các biểu tượng đã trở thành quy ước chung của nhân loại (Ngôi nhà = Trang chủ, Kính lúp = Tìm kiếm, Bánh răng = Cài đặt, Chiếc chuông = Thông báo).
\item \textbf{Tránh quá chi tiết (Simplicity):} Loại bỏ mọi chi tiết thừa rườm rà. Icon hiển thị trên điện thoại rất nhỏ, nếu quá nhiều chi tiết sẽ bị nhòe và biến dạng.
\item \textbf{Dễ nhận diện (Recognizability):} Người dùng nhìn lướt qua trong 0.1 giây là phải hiểu ngay ý nghĩa đại diện.
\end{enumerate}

\bigskip\hrule\bigskip

\section{3.10 Thiết kế Nút nhấn (Buttons)}

\begin{itemize}[leftmargin=1.5em]
\item Nút nhấn là thành phần kích hoạt hành động (\textbf{Call-to-Action}) quan trọng nhất trên giao diện.
\item \textbf{Các yếu tố thiết kế:}
\item \textit{Nhãn nút (Button Label):} Phải là một động từ hành động rõ ràng (ví dụ: \textit{"Tạo tài khoản"}, \textit{"Gửi đơn hàng"} thay vì ghi chung chung \textit{"Đồng ý"}, \textit{"OK"}).
\item \textit{Hình dạng và Bo góc:} Bo góc nhẹ ($4 - 8$ px) tạo cảm giác hiện đại, thân thiện.
\item \textit{Phân cấp nút bấm (Button Hierarchy):}
\item \textbf{Nút chính (Primary Button):} Nổi bật nhất trên màn hình, chỉ có 1 nút chính duy nhất tại một khung nhìn.
\item \textbf{Nút phụ (Secondary Button):} Màu trung tính hoặc dạng viền mảnh (Ghost Button).
\item \textbf{Nút hủy/nguy hiểm (Destructive Button):} Tông màu đỏ cảnh báo.
\item \textit{Kích thước cảm ứng (Touch Target):} Tối thiểu \textbf{$44 \times 44$ pixel} trên màn hình di động theo tiêu chuẩn công thái học để ngón tay cái người dùng không bấm trượt.
\end{itemize}

\bigskip\hrule\bigskip

\section{3.11 Thiết kế Biểu mẫu nhập liệu (Form Design)}

\begin{itemize}[leftmargin=1.5em]
\item Biểu mẫu là cầu nối quan trọng nhất để thu thập thông tin khách hàng.
\item \textbf{Nguyên tắc thiết kế biểu mẫu chuyên nghiệp:}
\item Bố trí các trường nhập liệu theo đúng trình tự thực tế và thói quen tâm lý của con người.
\item Bố trí theo cấu trúc 1 cột duy nhất (Single-column layout) giúp mắt người di chuyển từ trên xuống dưới mượt mà, tăng tỷ lệ hoàn thành form lên đến 30\% so với bố trí 2 cột zíc zắc.
\item Hỗ trợ hoàn hảo phím \textbf{Tab} để di chuyển tuần tự qua các ô.
\item \textbf{Cơ chế kiểm tra lỗi trực tiếp (Inline Validation):} Hiển thị dấu tích xanh hoặc thông báo đỏ ngay khi người dùng vừa rời khỏi ô nhập liệu, không đợi đến khi bấm nút Submit mới báo lỗi.
\end{itemize}

\bigskip\hrule\bigskip

\section{3.12 Thiết kế Cửa sổ con và Hộp thoại thông báo (Popups \& Dialogs)}

\begin{itemize}[leftmargin=1.5em]
\item Sử dụng cửa sổ con (Popup / Modal Dialog) khi cần thu hút sự tập trung tuyệt đối của người dùng vào một tác vụ khẩn cấp hoặc xác nhận quan trọng.
\item \textbf{Tin nhắn thông báo (Messages):} Ngắn gọn, súc tích, văn phong tôn trọng người dùng, nói rõ chuyện gì đã xảy ra và đưa ra nút giải quyết trực tiếp.
\item \textbf{Nhắc nhở hướng dẫn (Instructional Prompts):} Đặt đúng vị trí theo ngữ cảnh (Contextual Tooltip) và chỉ xuất hiện khi người dùng thực sự cần, không làm phiền trải nghiệm tự nhiên.
\end{itemize}

\bigskip\hrule\bigskip

\section{3.13 So sánh nền tảng: Website và Web Application}

\begin{table}[htbp]
\centering
\begin{tabularx}{\textwidth}{p{0.2\textwidth} p{0.38\textwidth} p{0.38\textwidth}}
\toprule
\textbf{Tiêu chí} & \textbf{Website Thông Thường} & \textbf{Web Application (Web App)} \\
\midrule
\textbf{Bản chất \& Mục đích} & Tập hợp các trang thông tin tĩnh/bán tĩnh nhằm cung cấp nội dung cho người đọc (Tin tức, Blog, Giới thiệu công ty). & Chương trình phần mềm tương tác chạy trên trình duyệt, thực hiện các tác vụ tính toán và xử lý dữ liệu phức tạp. \\
\textbf{Khả năng tương tác (Interactivity)} & Tương tác chủ yếu là một chiều: Đọc văn bản, xem hình ảnh, bấm liên kết đọc bài viết. & Tương tác hai chiều mạnh mẽ: Nhập dữ liệu, chỉnh sửa tài liệu thời gian thực, quản lý luồng nghiệp vụ. \\
\textbf{Khả năng tích hợp (Integration)} & Tích hợp cơ bản các tiện ích: Form liên hệ, mạng xã hội, Google Maps. & Tích hợp sâu rộng hệ thống API phức tạp: Cổng thanh toán, Cơ sở dữ liệu đám mây, Phân tích dữ liệu lớn. \\
\textbf{Khả năng xác thực (Authentication)} & Thường không yêu cầu đăng nhập; ai cũng có thể truy cập đọc nội dung tự do. & Bắt buộc phải có hệ thống xác thực người dùng chặt chẽ và phân quyền truy cập đa cấp (Admin, User). \\
\textbf{Phong cách thiết kế (Design Language)} & Tập trung vào nghệ thuật trình bày bài viết, hình ảnh thương hiệu, tối ưu hóa SEO. & Tập trung vào bảng điều khiển (Dashboard), hệ thống thanh công cụ, các bảng dữ liệu và luồng thao tác. \\
\textbf{Ví dụ điển hình} & Báo VnExpress, Dân Trí, Trang web giới thiệu trường đại học. & Gmail, Google Docs, Trello, Figma Web, Bảng quản trị Shopify. \\
\bottomrule
\end{tabularx}
\end{table}

\bigskip\hrule\bigskip

\section{3.14 So sánh thiết kế giao diện: Web vs Windows Form (Desktop) vs Mobile App}

\begin{table}[htbp]
\centering
\begin{tabularx}{\textwidth}{p{0.2\textwidth} p{0.25\textwidth} p{0.25\textwidth} p{0.25\textwidth}}
\toprule
\textbf{Tiêu chí so sánh} & \textbf{Giao diện Web (Web UI)} & \textbf{Giao diện Desktop (WinForms/.NET)} & \textbf{Giao diện Di động (Mobile App)} \\
\midrule
\textbf{1. Môi trường \& Kỹ thuật} & Chạy trên trình duyệt qua HTML, CSS, JavaScript; phụ thuộc vào chuẩn web và engine trình duyệt. & Chạy trực tiếp trên Hệ điều hành Windows (.NET Framework, C\#, WinForms, WPF); tận dụng tối đa tài nguyên phần cứng máy tính. & Chạy trực tiếp trên Hệ điều hành di động iOS/Android (Swift, Kotlin, Flutter, React Native). \\
\textbf{2. Kích thước màn hình} & Cực kỳ đa dạng; bắt buộc phải thiết kế đáp ứng (\textbf{Responsive Web Design}) co giãn từ 320px đến 4K. & Màn hình nằm ngang cố định, diện tích hiển thị rộng rãi, độ phân giải thường từ Full HD trở lên. & Màn hình dọc nhỏ gọn, không gian hiển thị hạn hẹp; mật độ điểm ảnh (PPI) cực kỳ cao. \\
\textbf{3. Phương thức tương tác} & Chuột và bàn phím: Thao tác nhấp, rê chuột (Hover), cuộn trang mượt mà. & Chuột và bàn phím độ chính xác tuyệt đối từng pixel; hỗ trợ phím tắt và tổ hợp phím phức tạp. & Cảm ứng bằng ngón tay: Chạm (Tap), vuốt (Swipe), kéo thả, kẹp thu phóng (Pinch-to-zoom), rung haptic. \\
\textbf{4. Hướng dẫn thiết kế (Guidelines)} & Tuân theo các quy chuẩn Web quốc tế (W3C), Bootstrap Grid, Tailwind UI. & Tuân thủ các quy tắc kinh điển của Microsoft: Windows UI / Fluent Design System. & Tuân thủ nghiêm ngặt chuẩn mực: Apple Human Interface Guidelines (HIG) hoặc Google Material Design. \\
\bottomrule
\end{tabularx}
\end{table}


\bigskip\hrule\bigskip

\chapter{CHƯƠNG 4: ĐÁNH GIÁ GIAO DIỆN NGƯỜI DÙNG (HEURISTIC EVALUATION)}

\section{4.1 Khái niệm Heuristic Evaluation}

\subsection{4.1.1 Heuristic Evaluation là gì?}
\begin{itemize}[leftmargin=1.5em]
\item \textbf{Heuristic Evaluation (Đánh giá giao diện theo bộ quy tắc kinh nghiệm):} Là một phương pháp kiểm định khả dụng (\textit{Usability Inspection Method}) phổ biến bậc nhất trong ngành công nghệ phần mềm.
\item Trong phương pháp này, một hoặc một nhóm các chuyên gia UI/UX sẽ trực tiếp soi chiếu, đánh giá và đo lường tính dễ-sử-dụng (\textit{Usability}) của giao diện người dùng dựa trên một tập hợp các quy tắc thiết kế đã được thực nghiệm và công nhận rộng rãi trên thế giới.
\end{itemize}

\bigskip\hrule\bigskip

\subsection{4.1.2 Nguồn gốc và Giá trị thực tiễn}
\begin{itemize}[leftmargin=1.5em]
\item Phương pháp được khởi xướng và hoàn thiện bởi Tiến sĩ \textbf{Jakob Nielsen} (đồng sáng lập tập đoàn tư vấn Nielsen Norman Group) và Rolf Molich vào những năm 1990, kết hợp cùng các công trình nghiên cứu của \textbf{Ben Shneiderman} trong cuốn sách kinh điển \textit{"Designing the User Interface"}.
\item \textbf{Giá trị thực tiễn:}  
\item Giúp phát hiện nhanh chóng từ \textbf{70\% đến 80\% các vấn đề về khả dụng} ngay từ giai đoạn bản vẽ thiết kế sơ bộ.
\item Tiết kiệm tối đa thời gian và chi phí so với việc tổ chức các phòng thí nghiệm kiểm thử người dùng thực tế tốn kém.
\end{itemize}

\begin{figure}[htbp]
\centering
\includegraphics[width=0.75\textwidth]{images/image54.png}
\caption{Phương pháp đánh giá giao diện Heuristic Evaluation}
\end{figure}

\bigskip\hrule\bigskip

\section{4.2 Mười nguyên tắc vàng Heuristic của Jakob Nielsen}

Mười nguyên tắc Heuristic kinh điển của Jakob Nielsen đã trở thành bộ luật chuẩn mực toàn cầu để đánh giá mọi giao diện số:


\begin{verbatim}
                            ┌───────────────────────────────────────────┐
                            │   10 NGUYÊN TẮC HEURISTIC JAKOB NIELSEN   │
                            └─────────────────────┬─────────────────────┘
                                                  │
       ┌──────────────┬──────────────┬────────────┴─┬──────────────┬──────────────┐
       ▼              ▼              ▼              ▼              ▼              ▼
1. Trạng thái  2. Thế giới    3. Kiểm soát   4. Nhất quán   5. Phòng chống 6. Gợi ý thay
   hệ thống       thực           tự do          chuẩn hóa       lỗi            vì nhớ
       │              │              │              │
       ▼              ▼              ▼              ▼
 7. Linh hoạt   8. Thẩm mỹ     9. Khắc phục   10. Trợ giúp &
    hiệu quả       tối giản       lỗi dễ dàng     tài liệu
\end{verbatim}

\bigskip\hrule\bigskip

\subsection{Nguyên tắc 1: Sự hiển thị của trạng thái hệ thống (Visibility of System Status)}
\begin{itemize}[leftmargin=1.5em]
\item \textbf{Nội dung:} Hệ thống phải luôn thông báo cho người dùng biết chuyện gì đang diễn ra, thông qua các phản hồi trực quan, thích hợp và kịp thời trong một khoảng thời gian chấp nhận được.
\item \textbf{Yêu cầu thiết kế:}  
\item Khi người dùng thực hiện một hành động (tải dữ liệu, gửi form, lưu file), hệ thống phải có tín hiệu phản hồi ngay tức thì (dưới 0.1 giây).
\item Đối với các tác vụ tốn thời gian, bắt buộc phải cung cấp thanh tiến trình (\textbf{Progress bar}), biểu tượng vòng xoay tải (\textbf{Spinner}), hoặc hiển thị thời gian ước tính còn lại.
\item \textbf{Ví dụ thực tế:} Hệ điều hành Apple macOS và iOS luôn hiển thị thanh tiến trình cài đặt kèm dòng chữ ước tính chính xác: \textit{"Còn lại khoảng 15 phút..."}. Trên ứng dụng Grab/ShopeeFood, người dùng theo dõi được vị trí của tài xế di chuyển thời gian thực trên bản đồ vệ tinh.
\end{itemize}

\begin{figure}[htbp]
\centering
\includegraphics[width=0.75\textwidth]{images/image55.png}
\caption{Nguyên tắc 1: Hiển thị trạng thái hệ thống thời gian thực}
\end{figure}

\bigskip\hrule\bigskip

\subsection{Nguyên tắc 2: Liên kết giữa hệ thống và thực tế (Match Between System \& the Real World)}
\begin{itemize}[leftmargin=1.5em]
\item \textbf{Nội dung:} Thiết kế giao diện cần sử dụng ngôn ngữ, khái niệm, hình ảnh và biểu tượng bắt nguồn từ cuộc sống đời thực quen thuộc của người dùng, thay vì dùng các thuật ngữ kỹ thuật mang tính chất lập trình trừu tượng.
\item \textbf{Yêu cầu thiết kế:}  
\item Các thông tin và chuỗi hành động phải được sắp xếp theo đúng trình tự tự nhiên và logic của thế giới thực.
\item \textbf{Ví dụ thực tế:}  
\item Chiếc biểu tượng thùng rác (\textit{Recycle Bin / Trash}) dùng để chứa các tệp tin đã xóa – hình ảnh thùng rác ngoài đời thực được chuyển hóa thành một ẩn dụ kỹ thuật số hoàn hảo.
\item Trên ứng dụng GrabFood, menu hiển thị hình ảnh chụp các đĩa thức ăn thật ngon mắt; trong khi ứng dụng BAEMIN sử dụng các hình vẽ minh họa đồ ăn dí dỏm, thân thuộc với phong cách ẩm thực đường phố Việt Nam.
\end{itemize}

\begin{figure}[htbp]
\centering
\includegraphics[width=0.75\textwidth]{images/image56.png}
\caption{Nguyên tắc 2: Tương thích hoàn hảo với thế giới thực}
\end{figure}

\bigskip\hrule\bigskip

\subsection{Nguyên tắc 3: Kiểm soát người dùng và sự tự do (User Control \& Freedom)}
\begin{itemize}[leftmargin=1.5em]
\item \textbf{Nội dung:} Người dùng thường xuyên thực hiện các thao tác sai lầm do bất cẩn, vội vàng hoặc hiểu nhầm. Hệ thống bắt buộc phải cung cấp một \textbf{"Cửa thoát hiểm khẩn cấp" (Emergency Exit)} được đánh dấu rõ ràng, cho phép người dùng thoát khỏi trạng thái không mong muốn ngay lập tức mà không phải trải qua các thủ tục phức tạp.
\item \textbf{Yêu cầu thiết kế:}  
\item Cung cấp các nút bấm: Hủy bỏ (\textit{Cancel}), Đóng (\textit{Close}), Thoát (\textit{Exit}).
\item Hỗ trợ hoàn hảo tính năng Hoàn tác (\textbf{Undo}) và Làm lại (\textbf{Redo}).
\item Luôn hỏi xác nhận trước khi thực hiện các hành động có tính chất phá hủy dữ liệu nghiêm trọng.
\item \textbf{Ví dụ thực tế:} Ứng dụng Gmail cung cấp nút "Hoàn tác" (Undo Send) xuất hiện ở góc màn hình trong vòng 30 giây ngay sau khi người dùng bấm gửi thư, cứu nguy cho hàng triệu trường hợp gửi nhầm email.
\end{itemize}

\begin{figure}[htbp]
\centering
\includegraphics[width=0.75\textwidth]{images/image57.png}
\caption{Nguyên tắc 3: Đảm bảo quyền kiểm soát và tự do cho người dùng}
\end{figure}

\bigskip\hrule\bigskip

\subsection{Nguyên tắc 4: Tính đồng bộ và chuẩn hóa (Consistency \& Standards)}
\begin{itemize}[leftmargin=1.5em]
\item \textbf{Nội dung:} Không bắt người dùng phải suy đoán xem những từ ngữ, tình huống hoặc hành động khác nhau có cùng mang một ý nghĩa hay không.
\item \textbf{Phân loại nhất quán:}
\item \textit{Nhất quán nội bộ (Internal Consistency):} Duy trì đồng nhất về bảng màu, vị trí nút bấm, kiểu chữ và cách dùng từ trên tất cả các màn hình bên trong cùng một sản phẩm.
\item \textit{Nhất quán bên ngoài (External Consistency):} Tuân thủ các quy ước chuẩn mực chung của toàn ngành hoặc của hệ điều hành nền tảng (ví dụ: Logo luôn nằm ở góc trên bên trái; biểu tượng bánh răng luôn là Cài đặt; màu đỏ là cảnh báo).
\item \textbf{Ví dụ thực tế:} Hệ sinh thái phần mềm đồ sộ của Microsoft (Word, Excel, PowerPoint, Outlook) duy trì sự đồng bộ chuẩn mực tuyệt đối từ hệ thống thanh công cụ Ribbon, tổ hợp phím tắt cho đến các biểu tượng thao tác.
\end{itemize}

\begin{figure}[htbp]
\centering
\includegraphics[width=0.75\textwidth]{images/image58.png}
\caption{Nguyên tắc 4: Tính đồng bộ và chuẩn hóa trong toàn hệ thống}
\end{figure}

\bigskip\hrule\bigskip

\subsection{Nguyên tắc 5: Ngăn chặn lỗi xảy ra (Error Prevention)}
\begin{itemize}[leftmargin=1.5em]
\item \textbf{Nội dung:} Một thiết kế cẩn trọng giúp \textbf{ngăn ngừa lỗi phát sinh ngay từ đầu} có giá trị cao hơn gấp nhiều lần việc cung cấp một thông báo lỗi hay nhất.
\item \textbf{Yêu cầu thiết kế:}  
\item Loại bỏ hoàn toàn các điều kiện dễ dẫn đến nhầm lẫn.
\item Cung cấp các ràng buộc thông minh: Ví dụ, ô nhập số điện thoại chỉ cho phép gõ số, không cho gõ chữ; ô chọn ngày sinh không cho phép chọn ngày trong tương lai.
\item Vô hiệu hóa nút bấm xác nhận (Disable submit button) cho đến khi người dùng điền đầy đủ và đúng định dạng các trường bắt buộc.
\item \textbf{Ví dụ thực tế:} Trên các trang web đặt vé máy bay hiện đại, khi hành khách chọn ngày bay, bộ chọn lịch (\textit{Datepicker}) sẽ tự động làm mờ và khóa toàn bộ các ngày về nhỏ hơn ngày đi, loại trừ hoàn toàn khả năng khách hàng đặt ngược ngày bay.
\end{itemize}

\begin{figure}[htbp]
\centering
\includegraphics[width=0.75\textwidth]{images/image59.png}
\caption{Nguyên tắc 5: Chủ động phòng chống và ngăn chặn lỗi}
\end{figure}

\bigskip\hrule\bigskip

\subsection{Nguyên tắc 6: Gợi ý tốt hơn là gợi nhớ (Recognition Rather Than Recall)}
\begin{itemize}[leftmargin=1.5em]
\item \textbf{Nội dung:} Giảm thiểu tối đa gánh nặng ghi nhớ cho bộ não con người bằng cách làm cho các đối tượng, hành động và tùy chọn luôn luôn hiển thị trực quan trước mắt.
\item \textbf{Quy luật tâm lý:} Con người nhận diện (Recognition) thông tin khi nhìn thấy dễ dàng hơn gấp hàng chục lần so với việc phải tự nhớ lại (Recall) thông tin từ trong ký ức trống rỗng.
\item \textbf{Ví dụ thực tế:}  
\item Hộp tìm kiếm của Google tự động xổ danh sách gợi ý từ khóa liên quan (\textbf{Google Autocomplete / Suggestions}) ngay khi người dùng mới gõ 1-2 ký tự đầu tiên.
\item Sàn thương mại điện tử Amazon, Shopee luôn dành riêng một khu vực hiển thị danh sách \textit{"Các sản phẩm bạn đã xem gần đây"}, giúp khách hàng dễ dàng tìm lại món hàng mà không cần phải nhớ tên sản phẩm.
\end{itemize}

\begin{figure}[htbp]
\centering
\includegraphics[width=0.75\textwidth]{images/image60.png}
\caption{Nguyên tắc 6: Gợi ý trực quan tốt hơn ép buộc nhớ lại}
\end{figure}

\bigskip\hrule\bigskip

\subsection{Nguyên tắc 7: Tính linh động và hiệu quả khi sử dụng (Flexibility \& Efficiency of Use)}
\begin{itemize}[leftmargin=1.5em]
\item \textbf{Nội dung:} Một thiết kế giao diện xuất sắc phải phục vụ hài hòa cho cả hai nhóm đối tượng: người mới làm quen (\textit{Novice users}) và người sử dụng chuyên nghiệp thành thạo (\textit{Expert users}).
\item \textbf{Yêu cầu thiết kế:}  
\item Cho phép người dùng tùy biến các thao tác thường dùng.
\item Cung cấp các công cụ tăng tốc (\textbf{Accelerators}), tổ hợp phím tắt nhanh và cử chỉ tiện ích.
\item \textbf{Ví dụ thực tế:}  
\item Người dùng phổ thông có thể dùng chuột bấm vào menu \textit{Edit $\rightarrow$ Copy $\rightarrow$ Paste}; trong khi người dùng thành thạo sử dụng ngay tổ hợp phím tắt \texttt{Ctrl + C} và \texttt{Ctrl + V}.
\item Tính năng chạm lực hoặc nhấn giữ (\textbf{3D Touch / Haptic Touch}) trên màn hình iPhone cho phép người dùng kích hoạt nhanh các lối tắt của ứng dụng mà không cần mở app.
\end{itemize}

\begin{figure}[htbp]
\centering
\includegraphics[width=0.75\textwidth]{images/image61.png}
\caption{Nguyên tắc 7: Tính linh động và hiệu quả cho mọi cấp độ người dùng}
\end{figure}

\bigskip\hrule\bigskip

\subsection{Nguyên tắc 8: Một thiết kế thẩm mỹ và tối giản (An Aesthetic \& Minimalist Design)}
\begin{itemize}[leftmargin=1.5em]
\item \textbf{Nội dung:} Giao diện không nên chứa đựng những thông tin không liên quan hoặc hiếm khi cần thiết.
\item \textbf{Quy luật cạnh tranh thị giác:} Mỗi một đơn vị thông tin rườm rà được nhồi nhét thêm vào màn hình sẽ cạnh tranh trực tiếp sự chú ý với các đơn vị thông tin quan trọng khác, và làm giảm khả năng nhìn thấy của các yếu tố cốt lõi.
\item \textbf{Triết lý tối giản:} Tối giản hóa không có nghĩa là cắt xén sơ sài, mà là nghệ thuật chắt lọc để \textbf{mỗi một chi tiết xuất hiện đều phải mang một giá trị nhất định}.
\item \textbf{Ví dụ thực tế:} Trang chủ của nền tảng đặt phòng toàn cầu \textbf{Airbnb} được tối giản hóa tới mức tối đa: menu được tinh gọn hoàn toàn, nhường toàn bộ sự tập trung cho một thanh tìm kiếm thanh lịch chỉ gồm 4 trường thông tin: \textit{Địa điểm – Ngày đến – Ngày đi – Số lượng khách}.
\end{itemize}

\begin{figure}[htbp]
\centering
\includegraphics[width=0.75\textwidth]{images/image62.png}
\caption{Nguyên tắc 8: Thẩm mỹ và thiết kế tối giản tinh tế}
\end{figure}

\bigskip\hrule\bigskip

\subsection{Nguyên tắc 9: Giúp người dùng nhận biết, chẩn đoán và khắc phục lỗi (Help Users Recognize, Diagnose, \& Recover from Errors)}
\begin{itemize}[leftmargin=1.5em]
\item \textbf{Nội dung:} Các thông báo lỗi phải được diễn đạt bằng ngôn từ bình dân, lịch sự, dễ hiểu; nói rõ chính xác nguyên nhân vấn đề và đề xuất hướng giải quyết mang tính xây dựng.
\item \textbf{Yêu cầu nghiêm ngặt:}  
\item Tuyệt đối không dùng những thông báo cộc lốc, vô trách nhiệm (ví dụ: \textit{"Lỗi không xác định!"}).
\item Tuyệt đối không quăng vào mặt người dùng phổ thông các mã lỗi lập trình bí hiểm (ví dụ: \textit{"NullPointerException at line 342"} hay \textit{"Error code: 0x80070002"}).
\item Nói rõ cho người dùng biết chuyện gì vừa xảy ra và họ cần phải làm gì tiếp theo để giải quyết.
\item \textbf{Ví dụ thực tế:}  
\item Hãng phim hoạt hình \textbf{PIXAR} sử dụng chính hình ảnh nhân vật Sadness trong bộ phim nổi tiếng \textit{Inside Out} trên trang báo lỗi \textbf{404 Not Found}, vừa xoa dịu tâm lý khó chịu của người dùng bằng sự hóm hỉnh, vừa cung cấp nút bấm đưa họ quay về trang chủ an toàn.
\item Hãng \textit{Email Center UK} thiết kế trang 404 bằng cách đưa ảnh chân dung của 4 nhân viên chịu trách nhiệm để người dùng chọn một người để phàn nàn, biến một trải nghiệm lỗi thành một khoảnh khắc thú vị và uy tín.
\end{itemize}

\bigskip\hrule\bigskip

\subsection{Nguyên tắc 10: Sự giúp đỡ và tài liệu hóa (Help \& Documentation)}
\begin{itemize}[leftmargin=1.5em]
\item \textbf{Nội dung:} Mặc dù lý tưởng nhất là một hệ thống được thiết kế hoàn hảo đến mức người dùng không cần bất kỳ sự giải thích nào vẫn thao tác mượt mà, nhưng việc cung cấp tài liệu hướng dẫn và hỗ trợ vẫn luôn là điều bắt buộc.
\item \textbf{Yêu cầu tài liệu:}  
\item Phải dễ dàng tìm kiếm tra cứu.
\item Tập trung vào các tác vụ cụ thể của người dùng.
\item Liệt kê các bước thực hiện ngắn gọn, súc tích kèm hình minh họa trực quan.
\item \textbf{Ví dụ thực tế:}  
\item Trung tâm hỗ trợ của \textbf{Apple} cung cấp mục lục tra cứu theo từng chủ đề cùng thanh tìm kiếm từ khóa thông minh, giúp người dùng giải đáp ngay thắc mắc mà không phải lội qua hàng trăm trang tài liệu.
\item Các sàn thương mại điện tử (Shopee, Tiki, Lazada) xây dựng trang Trung tâm trợ giúp với danh mục các câu hỏi thường gặp (\textbf{FAQ}) và hệ thống Chatbot trí tuệ nhân tạo giải đáp trực tuyến 24/7.
\end{itemize}

\begin{figure}[htbp]
\centering
\includegraphics[width=0.75\textwidth]{images/image63.png}
\caption{Nguyên tắc 10: Tài liệu hóa và hệ thống hỗ trợ người dùng toàn diện}
\end{figure}

\bigskip\hrule\bigskip

\section{4.3 Câu hỏi ôn tập trắc nghiệm \& câu hỏi ngắn củng cố kiến thức}

Dưới đây là các câu hỏi ôn tập trọng tâm giúp sinh viên củng cố vững chắc kiến thức:


\subsection{Câu 1:}
\begin{myquote}
\textbf{Câu hỏi:} Ứng dụng thương mại điện tử cung cấp tính năng tìm kiếm thông minh và tự động gợi ý từ khóa khi người dùng nhập vào ô tìm kiếm. Đây là ứng dụng của nguyên tắc Heuristic nào?
\textbf{Đáp án:} \textbf{Nguyên tắc số 6: Gợi ý tốt hơn là gợi nhớ (Recognition Rather Than Recall).}
\textit{Giải thích:} Tính năng này giúp người dùng nhận diện nhanh từ khóa họ cần tìm ngay trước mắt thay vì bắt bộ não phải nhớ lại chính xác từng ký tự của tên sản phẩm.
\end{myquote}

\subsection{Câu 2:}
\begin{myquote}
\textbf{Câu hỏi:} Khi người dùng nhập sai thông tin trong biểu mẫu đăng ký, hệ thống chỉ hiển thị một dòng thông báo vỏn vẹn: \textit{"Lỗi không xác định!"} mà không chỉ ra chính xác ô nào bị lỗi và cách khắc phục ra sao. Thiết kế này vi phạm nghiêm trọng nguyên tắc nào?
\textbf{Đáp án:} \textbf{Nguyên tắc số 9: Giúp người dùng nhận biết, chẩn đoán và khắc phục lỗi (Help Users Recognize, Diagnose, \& Recover from Errors).}
\textit{Giải thích:} Hệ thống không cho người dùng biết nguyên nhân gây lỗi và không đề xuất giải pháp, khiến người dùng hoàn toàn bất lực.
\end{myquote}

\subsection{Câu 3:}
\begin{myquote}
\textbf{Câu hỏi:} Hệ thống cho phép người dùng hoàn tác (Undo) hoặc làm lại (Redo) các thao tác vừa thực hiện. Đây là ví dụ tiêu biểu của nguyên tắc nào?
A. Hiệu suất và sự linh hoạt
B. Thẩm mỹ và thiết kế tối giản
C. Kiểm soát và tự do của người dùng
D. Nhất quán và tiêu chuẩn
\textbf{Đáp án:} \textbf{C. Kiểm soát và tự do của người dùng (User Control \& Freedom).}
\textit{Giải thích:} Cung cấp "lối thoát hiểm" giúp người dùng tự tin thao tác mà không sợ mắc phải sai lầm không thể cứu vãn.
\end{myquote}

\subsection{Câu 4:}
\begin{myquote}
\textbf{Câu hỏi:} Nguyên tắc "Hiển thị trạng thái hệ thống" (Visibility of System Status) có ý nghĩa cốt lõi là gì?
A. Hệ thống phải luôn luôn hiển thị rõ ràng trạng thái hiện tại và cung cấp phản hồi kịp thời cho người dùng.
B. Hệ thống phải ẩn đi mọi thông tin kỹ thuật để màn hình đẹp mắt.
C. Hệ thống chỉ cần phản hồi khi có sự cố xảy ra.
D. Người dùng không cần biết hệ thống đang làm gì sau hậu trường.
\textbf{Đáp án:} \textbf{A. Hệ thống phải luôn luôn hiển thị rõ ràng trạng thái hiện tại và cung cấp phản hồi kịp thời cho người dùng.}
\end{myquote}

\bigskip\hrule\bigskip

\section{4.4 Các câu hỏi tình huống thực tế (Scenario-Based Questions \& Solutions)}

\subsection{Tình huống 1: Phản hồi hệ thống (Visibility of System Status)}
\begin{itemize}[leftmargin=1.5em]
\item \textbf{Bối cảnh:} Bạn đang sử dụng một trang web thương mại điện tử để đặt mua hàng. Sau khi điền xong thông tin thẻ và nhấn nút \textit{"Thanh toán"}, màn hình hoàn toàn đứng im, không hiển thị bất kỳ thông báo, biểu tượng xoay hay hiệu ứng nào trong suốt 30 giây.
\item \textbf{Câu hỏi 1:} Tình huống này vi phạm nguyên tắc nào của Jakob Nielsen?
\item \textbf{Trả lời 1:} Vi phạm nghiêm trọng \textbf{Nguyên tắc số 1: Sự hiển thị của trạng thái hệ thống (Visibility of System Status)}. Khi người dùng thực hiện một hành động mang tính chất tài chính nhạy cảm như thanh toán, việc hệ thống im bặt sẽ gây ra tâm lý hoang mang tột độ, người dùng sợ bị trừ tiền nhiều lần hoặc nghĩ trang web đã bị treo.
\item \textbf{Câu hỏi 2:} Nếu bạn là UX Designer, bạn sẽ đề xuất giải pháp nào để cải thiện?
\item \textbf{Trả lời 2 (Giải pháp đề xuất):}
\end{itemize}
\begin{enumerate}[leftmargin=1.5em]
\item Ngay lập tức làm mờ nền và hiển thị lớp phủ (\textbf{Loading Overlay}) kèm vòng xoay tiến trình và dòng chữ: \textit{"Đang xử lý giao dịch an toàn, vui lòng không tắt trình duyệt..."}.
\item Vô hiệu hóa nút bấm thanh toán (Disable button) để ngăn khách hàng mất kiên nhẫn bấm liên tiếp nhiều lần gây trùng lặp đơn hàng.
\item Cung cấp thanh tiến trình hiển thị thời gian ước tính hoàn tất.
\item Nếu quá thời gian quy định (Timeout), phải hiển thị ngay thông báo lỗi rõ ràng kèm mã tra cứu và nút \textit{"Thử lại"} hoặc liên hệ tổng đài hỗ trợ.
\end{enumerate}

\bigskip\hrule\bigskip

\subsection{Tình huống 2: Kiểm soát và tự do của người dùng (User Control \& Freedom)}
\begin{itemize}[leftmargin=1.5em]
\item \textbf{Bối cảnh:} Một ứng dụng ngân hàng số không cho phép người dùng hủy giao dịch chuyển tiền ngay sau khi phát hiện vừa nhập sai số tài khoản của người nhận, dù chưa đến bước xác thực OTP cuối cùng.
\item \textbf{Câu hỏi 1:} Nguyên tắc nào đang bị vi phạm?
\item \textbf{Trả lời 1:} Vi phạm \textbf{Nguyên tắc số 3: Kiểm soát và tự do của người dùng (User Control \& Freedom)}. Hệ thống đã tước đoạt quyền tự quyết và lối thoát hiểm khẩn cấp của khách hàng.
\item \textbf{Câu hỏi 2:} Đề xuất một giải pháp UI/UX giúp cải thiện trải nghiệm người dùng?
\item \textbf{Trả lời 2 (Giải pháp đề xuất):}
\end{itemize}
\begin{enumerate}[leftmargin=1.5em]
\item Bắt buộc hiển thị một \textbf{Màn hình xác nhận thông tin (Confirmation Screen)} trung gian trước bước nhập mã OTP: hiển thị rõ ràng Tên người thụ hưởng viết hoa, Tên ngân hàng thụ hưởng và Số tiền chuyển.
\item Bố trí nút \textbf{"Hủy bỏ giao dịch"} hoặc \textbf{"Chỉnh sửa thông tin"} nổi bật bên cạnh nút "Tiếp tục", cho phép người dùng quay lại bước trước để sửa đổi số tài khoản dễ dàng.
\end{enumerate}

\bigskip\hrule\bigskip

\subsection{Tình huống 3: Nhất quán và tiêu chuẩn (Consistency \& Standards)}
\begin{itemize}[leftmargin=1.5em]
\item \textbf{Bối cảnh:} Một ứng dụng đặt xe công nghệ hiển thị nút \textit{"Đặt xe"} bằng màu xanh dương thương hiệu trên màn hình chính, nhưng khi chuyển sang màn hình xem chi tiết chuyến đi, nút \textit{"Đặt xe"} lại bất ngờ chuyển sang màu đỏ rực.
\item \textbf{Câu hỏi 1:} Vấn đề này ảnh hưởng thế nào đến trải nghiệm người dùng?
\item \textbf{Trả lời 1:} Vi phạm \textbf{Nguyên tắc số 4: Tính đồng bộ và chuẩn hóa (Consistency \& Standards)}. Màu đỏ theo quy ước tâm lý chung là màu của cảnh báo nguy hiểm hoặc nút Hủy chuyến. Việc đổi màu bất thường khiến người dùng bị hoang mang, khựng lại và sợ bấm nhầm vào nút hủy.
\item \textbf{Câu hỏi 2:} Đề xuất cách khắc phục để đảm bảo tính nhất quán?
\item \textbf{Trả lời 2 (Giải pháp đề xuất):}  
\end{itemize}
Quy chuẩn hóa toàn bộ hệ thống nút bấm trong Design System: Nút hành động chính (Primary Call-to-Action) như "Đặt xe" phải duy trì duy nhất một màu xanh dương đồng nhất trên 100\% các màn hình. Màu đỏ chỉ được dành riêng cho nút "Hủy chuyến đi".


\bigskip\hrule\bigskip

\subsection{Tình huống 4: Phòng tránh lỗi (Error Prevention)}
\begin{itemize}[leftmargin=1.5em]
\item \textbf{Bối cảnh:} Trên trang web đặt vé máy bay, nếu người dùng chọn ngày về nhỏ hơn ngày đi (ví dụ: ngày đi là 20/12 nhưng ngày về lại chọn 15/12), hệ thống vẫn cho phép bấm tiếp tục bình thường mà không hiển thị cảnh báo nào cho đến tận bước thanh toán cuối cùng mới báo lỗi.
\item \textbf{Câu hỏi 1:} Thiết kế này vi phạm nguyên tắc nào?
\item \textbf{Trả lời 1:} Vi phạm nghiêm trọng \textbf{Nguyên tắc số 5: Ngăn chặn lỗi xảy ra (Error Prevention)}.
\item \textbf{Câu hỏi 2:} Hãy đề xuất một giải pháp UI giúp người dùng phòng tránh lỗi này?
\item \textbf{Trả lời 2 (Giải pháp đề xuất):}  
\end{itemize}
Ứng dụng bộ chọn lịch thông minh (\textbf{Smart Datepicker}): Khi người dùng đã chọn ngày đi là 20/12, hệ thống sẽ tự động vô hiệu hóa (làm mờ và không thể click) tất cả các ngày từ 19/12 trở về trước trên lịch chọn ngày về. Người dùng chỉ có thể chọn ngày về bắt đầu từ ngày 20/12 trở đi.


\bigskip\hrule\bigskip

\subsection{Tình huống 5: Nhận diện hơn là nhớ lại (Recognition Rather Than Recall)}
\begin{itemize}[leftmargin=1.5em]
\item \textbf{Bối cảnh:} Một ứng dụng quản lý công việc và dự án nội bộ yêu cầu nhân viên phải tự tay gõ mã dự án bằng văn bản (ví dụ: \texttt{PRJ\_2026\_GUID\_009}) mỗi lần muốn thêm một nhiệm vụ mới, thay vì hiển thị danh sách các dự án có sẵn để chọn.
\item \textbf{Câu hỏi 1:} Tại sao việc bắt người dùng nhớ mã dự án là một vấn đề UX tồi tệ?
\item \textbf{Trả lời 1:} Vi phạm \textbf{Nguyên tắc số 6: Gợi ý tốt hơn là gợi nhớ (Recognition Rather Than Recall)}. Não người không được thiết kế để ghi nhớ các chuỗi ký tự khô khan, phức tạp. Việc bắt gõ tay gây lãng phí thời gian, tăng gánh nặng nhận thức và dẫn đến tỷ lệ gõ sai mã cực kỳ cao.
\item \textbf{Câu hỏi 2:} Làm thế nào để cải thiện chức năng này?
\item \textbf{Trả lời 2 (Giải pháp đề xuất):}  
\end{itemize}
Thay thế ô nhập văn bản đơn thuần bằng một \textbf{Hộp chọn tìm kiếm thông minh (Searchable Dropdown / Select with Autocomplete)}: Hiển thị danh sách tên dự án trực quan kèm logo dự án. Người dùng chỉ cần gõ 1-2 chữ cái của tên dự án là hệ thống tự động lọc và gợi ý để họ chọn nhanh chóng.


\bigskip\hrule\bigskip

\section{4.5 Bài tập thực hành tổng hợp và Đề tài thiết kế lớn}

\subsection{Bài tập 1: Phân tích lỗi giao diện thực tế (Heuristic Audit)}
\begin{itemize}[leftmargin=1.5em]
\item \textbf{Yêu cầu:} Lựa chọn một website hoặc ứng dụng di động phổ biến tại Việt Nam (Shopee, Tiki, Lazada, VnExpress...).
\item Tiến hành đóng vai chuyên gia Heuristic, chụp lại ít nhất 3 màn hình có vấn đề về khả dụng.
\item Chỉ rõ từng màn hình vi phạm nguyên tắc nào trong 10 nguyên tắc của Jakob Nielsen và đưa ra phương án thiết kế khắc phục.
\end{itemize}

\subsection{Bài tập 2: So sánh trải nghiệm đối đầu giữa hai ứng dụng cùng lĩnh vực}
\begin{itemize}[leftmargin=1.5em]
\item \textbf{Yêu cầu:} Chọn hai ứng dụng cạnh tranh trực tiếp (ví dụ: \textit{Facebook Messenger} vs \textit{Zalo}, hoặc \textit{Grab} vs \textit{Be}).
\item So sánh cách bố trí giao diện, luồng thao tác và sự tuân thủ các nguyên tắc Heuristic giữa hai sản phẩm.
\item Chỉ ra sản phẩm nào mang lại trải nghiệm vượt trội hơn và giải thích lý do cụ thể.
\end{itemize}

\subsection{Bài tập 3: Kiểm tra cơ chế nhận diện và phục hồi lỗi hệ thống}
\begin{itemize}[leftmargin=1.5em]
\item \textbf{Yêu cầu:} Truy cập một website bất kỳ và cố tình thực hiện các thao tác sai (nhập sai email, để trống trường bắt buộc, bấm vào đường link không tồn tại để vào trang 404).
\item Ghi nhận và đánh giá cách thức hệ thống phản hồi lỗi theo nguyên tắc số 5 và số 9.
\end{itemize}

\subsection{Bài tập 4: Nghiên cứu chuyên sâu UX của một siêu ứng dụng (TikTok / Facebook)}
\begin{itemize}[leftmargin=1.5em]
\item \textbf{Yêu cầu:} Phân tích cách thức TikTok hoặc Facebook ứng dụng các nguyên tắc:
\item \textit{Thao tác trực tiếp (Direct Manipulation):} Vuốt dọc chuyển video tức thì.
\item \textit{Tính tối giản (Minimalist Design):} Toàn bộ màn hình là nội dung video, các nút tương tác dồn sang mép phải.
\item \textit{Hiệu suất sử dụng:} Thuật toán gợi ý cá nhân hóa cao độ.
\end{itemize}

\bigskip\hrule\bigskip

\subsection{Đề tài thiết kế lớn: Xây dựng giao diện Website Thương mại điện tử xanh}
\begin{myquote}
\textbf{Chủ đề thiết kế:} \textbf{"Bán thực phẩm tươi sạch góp phần bảo vệ môi trường"}
\textbf{Nhiệm vụ của sinh viên:}
1. Xây dựng 2 User Personas đại diện cho nhóm khách hàng mua thực phẩm sạch (nội trợ hiện đại và nhân viên văn phòng bận rộn).
2. Thiết lập sơ đồ trang web (Sitemap) và luồng công việc (Workflow) mua hàng trực tuyến.
3. Thiết kế hệ thống Wireframe các trang chính: Trang chủ, Trang danh mục sản phẩm 2 cấp tối ưu, Trang chi tiết sản phẩm minh bạch nguồn gốc và Trang thanh toán (áp dụng bài học \textit{Chiếc nút bấm 300 triệu USD}).
4. Ứng dụng hệ thống màu sắc xanh lá cây (Green) chủ đạo, phân cấp Typography chuẩn mực và tuân thủ tuyệt đối 10 nguyên tắc Heuristic của Jakob Nielsen.
\end{myquote}

\end{document}
